\documentclass[12pt,a4paper]{article}


%%%%%%%%%%%%%%%%%%%%%%%%%%%%%%%%%%%%% PACKAGES %%%%%%%%%%%%%%%%%%%%%%%%%%%%%%%%%%%%%%%%
%---------------------------biblatex Stuff
%% \usepackage[style=reading,
%             abstract=true,
%             sorting=ynt]{biblatex}
% \bibliography{../../../Bibliographies/Library.bib}
% \addbibresource{../../../Bibliographies/Library.bib}
% \usepackage[american]{babel}
% \usepackage{csquotes}
%---------------------------Predefined Stuff (not needed)
% \usepackage[final,nousetoc,nouselof,nouselot,hylinks,notitlepage,full]{}

%---------------------------Pulls the page out a bit makes it look better (in my opinion)
% \usepackage{a4wide}

%---------------------------Removes paragraph indentation (not needed most of the time now)
% \usepackage{parskip}

%---------------------------Allows inclusion of graphics easily and configurably
\usepackage{graphics} 
\usepackage{graphicx} 

%---------------------------Provides ways to make nice looking tables
\usepackage{booktabs}

%---------------------------Allows you to rotate tables and figures
\usepackage{rotating}
\usepackage{rotfloat}

%---------------------------Allows shading of table cells
\usepackage{colortbl}
% Define a simple command to use at the start of a table row to make it have a shaded background
\newcommand{\gray}{\rowcolor[gray]{.9}}

%---------------------------???
% \usepackage{textcomp}

\usepackage{marvosym}
% \usepackage{ wasysym }

%---------------------------Provides commands to make subfigures (figures with (a), (b) and (c))
\usepackage{subfigure}

%---------------------------Makes references hyperlinks in PDF output
\usepackage[bookmarks=false,                      % show bookmarks bar?
            unicode=false,                        % non-Latin characters in Acrobat’s bookmarks
            pdftoolbar=true,                      % show Acrobat’s toolbar?
            pdfmenubar=true,                      % show Acrobat’s menu?
            pdffitwindow=true,                    % window fit to page when opened
            pdfstartview={FitH},                  % fits the width of the page to the window
            pdftitle={Notes on Lagrangian for Atmospheric Propagation},      % title
            pdfauthor={D.S. Cargill},             % author
            pdfsubject={Notes on Lagrangian for Atmospheric Propagation},   % subject of the document
            pdfcreator={D.S. Cargill},            % creator of the document
            pdfproducer={D.S. Cargill},           % producer of the document
            pdfkeywords={key1} {key2} {key3},     % list of keywords
            pdfnewwindow=true,                    % links in new window
            colorlinks=true,                      % false: boxed links; true: colored links
            linkcolor=green,                      % color of internal links
            citecolor=blue,                       % color of links to bibliography
            filecolor=red,                        % color of file links
            urlcolor=cyan                         % color of external links
]{hyperref}

% \usepackage{url} (uneeded)

%---------------------------Provides ways to include syntax-highlighted source code
% \usepackage{listings}
% \lstset{frame=single, basicstyle=\ttfamily}

%---------------------------Provides Harvard-style referencing
% \usepackage{natbib}
% \bibpunct{(}{)}{;}{a}{,}{,}

%---------------------------for label replacement
%\usepackage{psfrag}   

%---------------------------Pretty much all of the ams maths packages
\usepackage{amsmath,amsthm,amssymb,amsfonts}

%---------------------------Fonts
    % \renewcommand{\sfdefault}{phv}
%    \renewcommand{\familydefault}{\sfdefault}
%    \usepackage{arev}
%    \usepackage[T1]{fontenc}

%---------------------------Titles
\usepackage{titling}

%---------------------------si units
%\usepackage{siunitx}

%---------------------------Provides good access to colours
\usepackage{xcolor}
% \usepackage{color} (uneeded)

%---------------------------Provides the count of last page
\usepackage{lastpage}

%---------------------------Spacing
\usepackage{setspace}
\singlespacing
% \onehalfspacing
% \doublespacing

%---------------------------Geometry: easy margins
% \usepackage[a4paper]{geometry}
% \usepackage[margin=1in]{geometry}
% the usual page size settings:
\headheight      =  0.25in   % we normally don't have running heads
\headsep         =  0.25in   % we normally don't have running heads
\topmargin       = -0.5in    % 1 inch from top edge of paper
\oddsidemargin   = -0.25in   % 0.0 = 1 inch from left edge of paper
\evensidemargin  = -0.25in   % 0.0 = 1 inch from left edge of paper
\textheight      =  9.0in
\textwidth       =  7.0in
\footskip        =  0.4in


%---------------------------Allows fancy stuff in the page header
\usepackage{fancyhdr}
% \setlength{\headheight}{25.3pt}
\pagestyle{fancy}
% \renewcommand{\headrulewidth}{0pt} %% Removes the line under header
  \renewcommand{\headrulewidth}{1pt}
  \renewcommand{\footrulewidth}{1pt}
%---------------------------Allows printout of the page layout
\usepackage{layout}

%---------------------------Vastly improves the standard formatting of captions
%\usepackage[small]{caption2}
% \usepackage[font=small,format=plain,labelfont=bf,up,textfont=it,up]{caption}                                    
\usepackage[margin=10pt,font=small,labelfont=bf, labelsep=endash]{caption}


%---------------------------Easy access to the Lorem Iip�sum dummy text
%\usepackage{lipsum}


%---------------------------Draft Watermark
%\usepackage[printwatermark]{xwatermark}
%\newwatermark[allpages,color=red!10,angle=0,scale=6,xpos=0,ypos=00]{DRAFT}
% \newwatermark[allpages,color=red!20,angle=0,scale=5,xpos=0,ypos=50]{DRAFT}
% \newwatermark[allpages,color=red!20,angle=0,scale=5,xpos=0,ypos=0]{DRAFT}
% \newwatermark[allpages,color=red!20,angle=0,scale=5,xpos=0,ypos=-50]{DRAFT}


%%%%%%%%%%%%%%%%%%%%%%%%%%%%%%%%%%%%% NEWCOMMANDS %%%%%%%%%%%%%%%%%%%%%%%%%%%%%%%%%%%%%%%%

%---------------------------Simple command I defined to allow me to mark TODO items in red
\newcommand{\todo}[1]{\textbf{\textcolor{red}{#1}}}


%---------------------------Simple command I defined to allow me to mark notes in blue
\newcommand{\anote}[1]{\textbf{\textcolor{blue}{#1}}}


% \DeclareCiteCommand{\wcite}
%     [\mkbibfootnote]
%     {}
%     {\printfield{title}--\printfield{url}}
%     {\addcomma\addspace}
%     {}
%%%%%%%%%%%%%%%%%%%%%%%%%%%%%%%%%%%%% INCLUSIONS %%%%%%%%%%%%%%%%%%%%%%%%%%%%%%%%%%%%%%%%
%---------------------------Macros
\def\boxit#1{\vbox{\hrule\hbox{\vrule\kern3pt
      \vbox{\kern3pt#1\kern3pt}\kern3pt\vrule}\hrule}}
\def\boxitd#1{\vspace{0.1in} \vbox{\hrule\hbox{\vrule\kern3pt
      \vbox{\abovedisplayskip=10pt \belowdisplayskip=10pt #1}
       \vrule}\hrule}}
\def\ind{\hspace{0.25in}}
\def\adb{\allowdisplaybreak}
\def\spb{\smallskip\par}
\def\mpb{\medskip\par}
\def\bpb{\bigskip\par}
\def\qed{{\ \vbox{\hrule\hbox{%
   \vrule height1.3ex\hskip0.8ex\vrule}\hrule
  }}\par}

\def\bm#1{\mbox{\boldmath{$#1$}}}
\def\tns#1{\bm{#1}}
\def\IN{\quad\mbox{in }}
\def\ON{\quad\mbox{on }}
\def\lb{\lambda}
\def\lbi{\lambda^{-1}}
\def\ep{\varepsilon}
\def\sg{\sigma}
\def\til#1{\tilde{#1}}
\def\tep{\bm{\ep}}
\def\tsg{\bm{\sg}}
\def\ttau{\bm{\tau}}
\def\tb{\bm{\beta}}
\def\tg{\bm{\gamma}}
\def\td{\bm{\delta}}
\def\enn{\frac{E}{1-\nu^2}}
\def\enon{\frac{E}{(1+\nu)(1-\ol{\nu})}}
\def\en{\frac{E}{1+\nu}}
\def\ol{\overline}
\def\ul{\underline}
\def\Ga{\Gamma}
\def\vs{\vspace{1\jot}}
\def\veta{\bm{\eta}}
\def\oveta{\ol{\bm{\eta}}}
\def\oth{\ol{\theta}}
\def\vphi{\bm{\phi}}
\def\vPhi{\bm{\Phi}}
\def\ovPhi{\ol{\bm{\Phi}}}
\def\vpsi{\bm{\psi}}
\def\ovphi{\ol{\bm{\phi}}}
\def\ovpsi{\ol{\bm{\psi}}}
\def\oovphi{\ol{\ol{\bm{\phi}}}}
\def\oovpsi{\ol{\ol{\bm{\psi}}}}
\def\olb{\ol{\lambda}}
\def\oolb{\ol{\olb}}
\def\vga{\bm{\gamma}}
\def\tz{\tilde z}
\def\Om{\Omega}
\def\oll{\ol{l}}
\def\om{\omega}
\def\tom{\tilde\om}
\def\oom{\ol{\omega}}
\def\oW{\ol{W}}
\def\oV{\ol{V}}
\def\ooV{\ol{\oV}}
\def\oQ{\ol{Q}}
\def\oP{\ol{P}}
\def\ov{\ol{v}}
\def\tu{\tilde{u}}
\def\hu{\hat{u}}
\def\hbeta{\hat{\beta}}
\def\toW{\widetilde{\oW}}
\def\hoW{\widehat{\oW}}
\def\tW{\widetilde{W}}
\def\ooom{\ol{\oom}}
\def\omu{\ol{\mu}}
\def\oomu{\ol{\omu}}
\def\ophi{\ol{\phi}}
\def\oPhi{\ol{\Phi}}
\def\opsi{\ol{\psi}}
\def\oophi{\ol{\ol{\phi}}}
\def\oopsi{\ol{\ol{\psi}}}
\def\oW{\ol{W}}
\def\ooW{\ol{\oW}}
\def\oR{\ol{r}}
\def\op{\ol{p}}
\def\os{\ol{s}}
\def\oq{\ol{q}}
\def\nn{\nonumber}
%\def\implies{\Rightarrow}
\def\nohyphens{\pretolerance 10000}
\def\bv{\bm{v}}
\def\bw{\bm{w}}
\def\bx{\bm{x}}
\def\dbx{\dot{\bx}}
\def\ddbx{\ddot{\bx}}
\def\dx{\dot{x}}
\def\ddx{\ddot{x}}
\def\dy{\dot{y}}
\def\ddy{\ddot{y}}
\def\df{\dot{f}}
\def\ddf{\ddot{f}}
\def\by{\bm{y}}
\def\bz{\bm{z}}
\def\bb{\bm{b}}
\def\bp{\bm{p}}
\def\bq{\bm{q}}
\def\be{\bm{e}}
\def\br{\bm{r}}
\def\bd{\bm{d}}
\def\bk{\bm{k}}
\def\bff{\bm{f}}
\def\bg{\bm{g}}
\def\bh{\bm{h}}
\def\bn{\bm{n}}
\def\bC{\bm{C}}
\def\bN{\bm{N}}
\def\bL{\bm{L}}
\def\bK{\bm{K}}
\def\bF{\bm{F}}
\def\bG{\bm{G}}
\def\bv{\bm{v}}
\def\bZ{\bm{Z}}
\def\bX{\bm{X}}
\def\bY{\bm{Y}}
\def\bS{\bm{S}}
\def\actf{\mathaccent{f}}
\def\Diamond{\diamond}
\def\b0{\mbox{\bf 0}}
\def\oh{\overline h}
\def\hbx{\hat{\bx}}
\def\Axb{$A \bx = \bb$}
\def\sAxb#1{$A^{(#1)} \bx = {\bb}^{(#1)}$}
\def\tAxb{$\tilde A \bx = \tilde {\bb}$}

%% Operator names
\def\tr{\mbox{tr}\,}
\def\vcurl{\mbox{\bf curl}\,}
\def\sdiv{\mbox{\rm div}\,}
\def\srot{\mbox{\rm rot}\,}
\def\vdiv{\mbox{\bf div}\,}
\def\vgrad{\mbox{\bf grad}\,}
\def\vgradh{\mbox{\bf grad}_h\,}
\def\intl{\int\limits}
\def\RWEOp#1{\nabla^2 #1 + k^2 #1 }
\def\WEOp#1{\nabla^2 #1 - \frac{1}{c^2} \frac{\partial^2}{\partial t^2} #1} 
\newcommand\conv[2]{\left[#1 \ast #2\right]}


\def\dim#1{\tilde{#1}}

%% Functions
\def\sinc#1{\frac{\mbox{\bf{sin}}{\left( #1 \right)}}{\left( #1 \right)}}
\def\sin#1{\mbox{\bf{sin}}{\left( #1 \right)}}
\def\cos#1{\mbox{\bf{cos}}{\left( #1 \right)}}
\def\tan#1{\mbox{\bf{tan}}{\left( #1 \right)}}
\def\cot#1{\mbox{\bf{cot}}{\left( #1 \right)}}
\def\sec#1{\mbox{\bf{sec}}{\left( #1 \right)}}
\def\csc#1{\mbox{\bf{csc}}{\left( #1 \right)}}
\def\erfc#1{\mbox{\bf{erfc}}\left( #1 \right)}
\def\erf#1{\mbox{\bf{erf}}\left( #1 \right)}
\def\ln#1{\mbox{\bf{ln}}\left( #1 \right)}
\def\sech#1{\mbox{\bf{sech}}{\left( #1 \right)}}
\def\cosh#1{\mbox{\bf{cosh}}{\left( #1 \right)}}
\def\tanh#1{\mbox{\bf{tanh}}{\left( #1 \right)}}



% theorems

% \newtheorem{theorem}{Theorem}[chapter]
% \newtheorem{corollary}{Corollary}[chapter]
% \newtheorem{lemma}{Lemma}
% \newenvironment{proof}{\noindent {\sc Proof}:}{\par \spb}
 %\newenvironment{example}{\noindent {\sc Example}:}{\par \spb}
 \newenvironment{define}{\noindent {\bf Definition}:}{\par \spb}

%Copyrights

% \newcommand\copyrightp{%
%         \null\vfil
%         \typeout{Copyright}
%         {\centering\bfseries
%                 \large\copyright\ \@year\\
%                 \@author\\
%                 \large ALL RIGHTS RESERVED\par}
%         \vfil\newpage}

% 

%%%%%%%%%%%%%%%%% Math Text %%%%%%%%%%%%%%%%%%%%%%%
\def\Des#1#2{#1_{\mathrm{#2}}}
% \newcommand{\Lnls}{L_{\mathrm{nls}}}
% \newcommand{\usol}{u_{\mathrm{sol}}}
% \newcommand{\Nnls}{N_{\mathrm{nls}}}

%%%%%%%%%%%%%%%%% Text %%%%%%%%%%%%%%%%%%%%%%%%%

\def\eg{e.\/g.\/, }
\def\ie{i.\/e.\/, }
\def\et{et al.\/, }
% \def\eg{e.\/g.\@\xspace }
% \def\ie{i.\/e.\@\xspace }
% \def\et{et al.\@\xspace }

%%%%%%%%%%%%%%%%% Statistics %%%%%%%%%%%%%%%%%%%%%%%%%

\def\Expect#1{\mathbb{E}\hspace{-0.03in}\left[ #1 \right]}
\def\Prob#1{\mathbb{P}\hspace{-0.03in}\left[ #1 \right]}
\def\Var#1{\mathbb{V}\hspace{-0.03in}\left[ #1 \right]}
\def\CV#1{\mathbb{C}_{\mathrm{var}}\hspace{-0.03in}\left[ #1 \right]}
\def\cor#1{\mathrm{cor}\hspace{-0.03in}\left(#1 \right)}
\def\cov#1{\mathrm{cov}\hspace{-0.03in}\left(#1 \right)}

%%%%%%%%%%%%%%%%% Functions %%%%%%%%%%%%%%%%%%%%%%%%%%

\def\cosh#1{\,\mathrm{cosh}\hspace{-0.03in}\left( #1 \right)}
\def\sinh#1{\,\mathrm{sinh}\hspace{-0.03in}\left( #1 \right)}
\def\tanh#1{\,\mathrm{tanh}\hspace{-0.03in}\left( #1 \right)}
\def\sech#1{\,\mathrm{sech}\hspace{-0.03in}\left( #1 \right)}
\def\csch#1{\,\mathrm{csch}\hspace{-0.03in}\left( #1 \right)}

\def\atanh#1{\,\mathrm{arctanh}\hspace{-0.03in}\left( #1 \right)}

\def\cos#1{\,\mathrm{cos}\hspace{-0.03in}\left( #1 \right)}
\def\sin#1{\,\mathrm{sin}\hspace{-0.03in}\left( #1 \right)}
\def\tan#1{\,\mathrm{tan}\hspace{-0.03in}\left( #1 \right)}
\def\sec#1{\,\mathrm{sec}\hspace{-0.03in}\left( #1 \right)}
\def\csc#1{\,\mathrm{csc}\hspace{-0.03in}\left( #1 \right)}

\def\atan#1{\,\mathrm{arctan}\hspace{-0.03in}\left( #1 \right)}
\def\acos#1{\mbox{\bf{arccos}}{\left( #1 \right)}}

\def\sgn#1{\,\mathrm{sgn}\hspace{-0.03in}\left( #1 \right)}


\def\heav#1{\,\mathrm{H}\hspace{-0.03in}\left( #1 \right)}

\def\ln#1{\,\mathrm{ln}\hspace{-0.03in}\left( #1 \right)}
% \def\exp#1{\,\mathrm{e}^{#1 }}
\def\exp#1{\,\mathrm{exp}\hspace{-0.03in}\left(#1 \right)}

\def\sinc#1{\,\mathrm{sinc}\hspace{-0.03in}\left( #1 \right)}

\def\erfc#1{\mbox{Erfc}\hspace{-0.03in}\left( #1 \right)}
\def\erf#1{\mbox{Erf}\hspace{-0.03in}\left( #1 \right)}

\def\For#1{\mathcal{F}\hspace{-0.03in}\left[ #1 \right]}
\def\iFor#1{\mathcal{F}^{-1}\hspace{-0.03in}\left[ #1 \right]}

\def\del#1{\,\delta\hspace{-0.03in}\left( #1 \right)}

\def\pdfNorm#1#2{\,\mathcal{N}\hspace{-0.03in}\left( #1,#2 \right)}
\def\Chi2#1{\,\Chi\hspace{-0.03in}\left( #1 \right)}

\def\circ#1{\,\mathrm{circ}\hspace{-0.03in}\left( #1 \right)}


%%%%%%%%%%%%%%%%% Derivatives %%%%%%%%%%%%%%%%%%%%%%%%%%

\newcommand\der[2]{\frac{d{#1}}{d{#2}}}
\newcommand\dder[2]{\frac{d^2{#1}}{d{#2}^2}}
\newcommand\ddder[2]{\frac{d^3{#1}}{d{#2}^3}}
\newcommand\pder[2]{\frac{\partial #1}{\partial #2}}
\newcommand\ppder[2]{\frac{\partial^2 #1}{\partial #2^2}}
\newcommand\pderpder[3]{\frac{\partial^2 #1}{\partial #2 \partial #3}}


\let\grad=\nabla

\let\truesum=\sum
\def\sum{\mathop{\textstyle\truesum}\limits}

\def\Re#1{\mathcal{\rm{Re}}\left[#1\right]}
\def\Im#1{\mathcal{\rm{Im}}\left[#1\right]}

\def\bigOh#1{O\hspace{-0.03in}\left(#1\right)}
\def\littelOh#1{o\hspace{-0.03in}\left(#1\right)}

% \let\<=\left<
% \let\>=\right>
\def\_#1{{\underline{#1}}}
\def\'#1{{\mathbf{#1}}}
\let\truehat=\^
\let\^=\hat

%\def\in{\mathrm{in}}
%\def\out{\mathrm{out}}
\def\tot{\mathrm{tot}}
\def\O#1{{\mbox{O}(#1)}}
\def\ufilt{u_{\mathrm{filt}}}
\def\utfilt{\tilde{u}_{\mathrm{filt}}}

\def\conj#1{\bar{#1}}
% \def\iint{\int \int}
\newcommand{\work}[2]{\ifnum\full=1{#1}\else{#2}\fi}


%%%%%%%%%%%%%%%%% Linear Algebra %%%%%%%%%%%%%%%%%%%%
\def\det#1{\,\mathrm{det}\hspace{-0.03in}\left( #1 \right)}



%%%%%%%%%%%%%%%%% Footnotes %%%%%%%%%%%%%%%%%%%%%%%%%

\renewcommand{\thefootnote}{\arabic{footnote}}   % Arabic numerals, e.g., 1, 2, 3...
% \renewcommand{\thefootnote}{\roman{footnote}}    % Roman numerals (lowercase), e.g., i, ii, iii...
% \renewcommand{\thefootnote}{\Roman{footnote}}    % Roman numerals (uppercase), e.g., I, II, III...
% \renewcommand{\thefootnote}{\alph{footnote}}     % Alphabetic (lowercase), e.g., a, b, c...
% \renewcommand{\thefootnote}{\Alph{footnote}}     % Alphabetic (uppercase), e.g., A, B, C...
% \renewcommand{\thefootnote}{\fnsymbol{footnote}} % A sequence of nine symbols (try it and see!)



% \newcommand{\intl}[4]{\int_{#3}^{#4} #1 \, d #2}
% \def\Del#1{\Delta(#1)}
\def\Int#1#2{\left[\int #1 \, d#2\right]}

\def\dmin#1#2{\displaystyle{\min_{#1}}\hspace{-0.03in}\left(#2 \right)}
\def\dmax#1#2{\displaystyle{\max_{#1}}\hspace{-0.03in}\left(#2 \right)}
\def\dargmin#1#2{\displaystyle{{\mathrm{argmin}}_{#1}}\hspace{-0.03in}\left(#2 \right)}
\def\dargmax#1#2{\displaystyle{{\mathrm{argmax}}_{#1}}\hspace{-0.03in}\left(#2 \right)}

\def\vect#1{\bm{#1}}
% \def\vect#1{\overrightarrow{#1}}
\def\hatvect#1{\vect{\hat{#1}}}
\def\gvect#1{\boldsymbol{#1}}
\def\matx#1{\bm{#1}}

% \def\red#1{\textcolor{red}{#1}}
% \def\blue#1{\textcolor{blue}{#1}}
% \def\green#1{\textcolor{green}{#1}}
% \def\cyan#1{\textcolor{cyan}{#1}}
% \def\magenta#1{\textcolor{magenta}{#1}}

\def\red#1{\color{red} #1 \color{black}}
\def\blue#1{\color{blue} #1 \color{black}}
\def\green#1{\color{green} #1 \color{black}}

\def\ip#1#2{\left \langle #1,#2 \right \rangle}
\def\Ip#1#2{\Re{\int \conj{#1}\, #2 \,dt}}
% \def\avg#1{\breve{#1}}
\def\avg#1{{#1}_{avg}}
\def\loc#1{{#1}_{loc}}
\def\var#1{{#1}_{var}}

\def\adjoint#1{#1^{\dagger}}
\def\Ep#1{\mathrm{E}\left[#1\right]}

\newenvironment{note}{\nopagebreak \begin{quote}
\noindent{\textcolor{red}{\textbf{DSC Note:}}}}{
\end{quote} \medskip} 




\def\iF#1{\mathcal{F}^{-1}\left[ #1 \right]}

\renewcommand{\thefootnote}{\fnsymbol{footnote}}

%%%%%%%%%%%%%%%%%%%%%%%%%%%%%%%%%%%%% TITLE STUFF %%%%%%%%%%%%%%%%%%%%%%%%%%%%%%%%%%%%%%%%

%----------------------------Title, Header and Footer
% Standard title, author etc.
\title{}
\author{Daniel S. Cargill, Ph.D.}
\date{\today}
% \affil{}

% Put text on the left-hand and right-hand side of the header
% \fancyhf{}
\lhead{Notes on Lagrangian for Atmospheric Propagation}
\chead{}
\rhead{\thedate}
\lfoot{\theauthor}
\cfoot{}
\rfoot{\thepage\ of \pageref{LastPage}}

\begin{document}

% \title{\LaTeX\ Template for Writing Papers}
% \author{Daniel S. Cargill}
\title{Notes: Derivation of ODEs for Atmospheric Propagation Using Lagrangian Approach}
\author{Daniel S. Cargill}
% \address{}
% \email{$^*$}{dc26@njit.edu}
% \emailandwebpage{}{dc26@njit.edu}{http://web.njit.edu/~dc26/}

\maketitle

\begin{abstract}


\end{abstract}

\section{Introduction}\label{sec:Introduction_and_Background}


\subsection{Helmholtz Equation for  Atmospheric Propagation}

A monochromatic EM source $E(\dim{x},\dim{y},\dim{z}) = A(\dim{x},\dim{y},\dim{z})\exp{-i\omega \dim{t}} + c.c.$, 
propagating in a nondispersive static atmosphere follows a stochastic Helmholtz equation

\begin{equation}
\begin{split}
  &\nabla^2 \dim{A} + k^2  (n_0 +  n(\dim{x},\dim{y},\dim{z}))^2 \dim{A} = 0 \quad  \mbox{for}   
  \quad -\infty < \dim{x},\dim{y} < \infty, \quad 0 < \dim{z} < L, \\
  &\quad  \mbox{where} \quad \dim{A}(\dim{x},\dim{y},0) = \dim{A}_{i}(\dim{x},\dim{y})\exp{i \Phi(\dim{x},\dim{y})},
\end{split}
\end{equation}
and $\dim{A}_{i}(\dim{x},\dim{y})$ is of finite support, i.e. $\dim{A}_{i}(\dim{x},\dim{y}) = 0$ 
when $\sqrt{\dim{x}^2 + \dim{y}^2} > D_{i}$. Note, $k = \frac{\omega}{c} = \frac{2\pi}{\lambda}$.  




%%%%%%%%%%%%%%%%%%%%%%%%%%%%%%%%%%%%%%%%%%%%%%%%%%%%%%%%%%%%%%%%%%%%%%%%%%%%%%%%%%%%%%%%%%%%%%%%%%%

\subsection{Atmospheric Statistics}
 Normally, the stochastic part of the atmosphere is described by the structure function~\cite{andrews2005laser}
\begin{equation}
\begin{split}
  \Expect{(n(\vect{\dim{r}}_1) - n(\vect{\dim{r}}_2))^2} &= 
  C_n^2  \begin{cases}
        l_{0}^{-4/3} \left|\vect{\dim{r}}_1-\vect{\dim{r}}_2\right|^{2} & 0 < \left|\vect{\dim{r}}_1-\vect{\dim{r}}_2\right| \leq l_{0} \\
       \left|\vect{\dim{r}}_1 - \vect{\dim{r}}_2 \right|^{2/3} &  l_{0}  \leq \left|\vect{\dim{r}}_1 - \vect{\dim{r}}_2\right| \leq  L_{0}  \\
   \end{cases}\\
   &= \Expect{n^2(\vect{\dim{r}}_1)} +  \Expect{n^2(\vect{\dim{r}}_2)} 
   - 2\Expect{n(\vect{\dim{r}}_1)n(\vect{\dim{r}}_2)} \\
   &= 2\sigma_n^2(1-\cor{n_1,n_2})
\end{split}
\end{equation}
so
\begin{equation}
  \cor{n_1,n_2} = 1 - 
    \begin{cases}
        \frac{C_n^2}{2\sigma_n^2}l_{0}^{-4/3} \left|\vect{\dim{r}}_1-\vect{\dim{r}}_{2}\right|^{2} & 0 < \left|\vect{\dim{r}}_1-\vect{\dim{r}}_{2}\right| \leq l_{0} \\
        \frac{C_n^2}{2\sigma_n^2} \left|\vect{\dim{r}}_1-\vect{\dim{r}}_{2}\right|^{2/3} &  l_{0}  \leq \left|\vect{\dim{r}}_1-\vect{\dim{r}}_{2}\right| \leq  L_{0}  \\
   \end{cases}
\end{equation}

If we assume that $|\cor{n_1,n_2}| \ll 1$ for $ \left|\vect{\dim{r}}_1-\vect{\dim{r}}_{2}\right| \gg L_0$, 
then letting $\left|\vect{\dim{r}}_1-\vect{\dim{r}}_{2}\right| = C_0 L_0$ for $C_0 > 1$ we can write
\begin{equation}
\begin{split}
  0 = 1 - C_0^{2/3} \frac{C_n^2 L_0^{2/3}}{2\sigma_n^2},
\end{split}
\end{equation}
or 
\begin{equation}
\begin{split}
  \sigma_n^2 =  C_0^{2/3} \frac{C_n^2 L_0^{2/3}}{2} 
\end{split}
\end{equation}
and
\begin{equation}
  \cor{n_1,n_2} = 1 - 
    \begin{cases}
        \frac{1}{C_0^{2/3}(L_0l_{0}^2)^{2/3}}\left|\vect{\dim{r}}_1-\vect{\dim{r}}_{2}\right|^{2} & 0 < \left|\vect{\dim{r}}_1-\vect{\dim{r}}_{2}\right| \leq l_{0} \\
        \frac{1}{C_0^{2/3}L_0^{2/3}} \left|\vect{\dim{r}}_1-\vect{\dim{r}}_{2}\right|^{2/3} &  l_{0}  \leq \left|\vect{\dim{r}}_1-\vect{\dim{r}}_{2}\right| \leq  L_{0}  \\
   \end{cases}
\end{equation}

\begin{table}[htp]
\begin{center}
\begin{tabular}{|c|c|c|}
\hline
Quantity                 & Symbol                    & Characteristic Value(s)  \\       
\hline
\hline
Wavelength               & $\lambda $                & $10^{-6} $ (meters)      \\                   
\hline
Wavenumber               & $k= \frac{2\pi}{\lambda}$ & $10^{6} $ (meters$^{-1}$)\\            
\hline
Kolmogorov Inner-scale   & $l_0$                     & $10^{-2} \leftrightarrow 10^{-3}$ (meters)       \\             
\hline
Kolmogorov Outer-scale   & $L_0$                     & $10 \leftrightarrow  10^{3}$ (meters)        \\            
\hline
Index Structure Constant & $C_n^2$                   & $10^{-20}\leftrightarrow 10^{-12}$ (meters$^{-2/3}$) \\
\hline
Index Variance           & $\sigma_n^2 =  C_0^{2/3} \frac{C_n^2 L_0^{2/3}}{2} $              & $C_0^{2/3} (10^{-18}\leftrightarrow 10^{-10})$ \\     
\hline
Apperture Diameter       & $D_{i}$                   & $1\leftrightarrow 10^{-2}$ (meters)     \\            
\hline
Propagation Distance     & $L$                       & $10^{3}\leftrightarrow 10^{4}$ (meters) \\     
\hline     
\end{tabular}
\end{center}
\caption{Atmospheric Scales}
\label{Atmospheric Scales}
\end{table}


%%%%%%%%%%%%%%%%%%%%%%%%%%%%%%%%%%%%%%%%%%%%%%%%%%%%%%%%%%%%%%%%%%%%%%%%%%%%%%%
\subsubsection{Nondimensionalization}

\begin{table}[htp]
\begin{center}
\begin{tabular}{|c|c|c|}
\hline
Quantity                 & Scale                & New Variable        \\       
\hline
\hline
$\dim{z}$                & $L_0$                & $\dim{z} = L_0 z$   \\                   
\hline
$\dim{x}$                & $l_0$                & $\dim{x} = l_0  x$  \\
\hline
$\dim{y}$                & $l_0$                & $\dim{y} = l_0  y$  \\
\hline
$n(l_0 \dim{x}, l_0 \dim{y}, L_0 \dim{z})$ & $\sigma_n$ & $n(l_0 \dim{x}, l_0 \dim{y}, L_0 \dim{z}) = \sigma_n n(x,y,z)$  \\
\hline
$\dim{A}(l_0 \dim{x}, l_0 \dim{y}, L_0 \dim{z})$ &  $A_0$  & $\dim{A}(l_0 \dim{x}, l_0 \dim{y}, L_0 \dim{z}) = A_0 A(x,y,z)$ \\
\hline   
\end{tabular}
\end{center}
\caption{Dimensional Scales}
\label{Dimensional Scales}
\end{table}


\begin{equation}
\begin{split}
  & \frac{l_0^2}{L_0^2} \ppder{A}{z} + \nabla_{\perp}^2 A + l_0^2 k^2 n_0^2 A 
  +  2 l_0^2k^2 n_0 \sigma_n n(x,y,z) A + l_0^2k^2 \sigma_n^2 n^2(x,y,z) A = 0 \\
  &\quad  \mbox{for}   
  \quad -\infty < x,y < \infty, \quad 0 < z < \frac{L}{L_0} \quad  
  \mbox{where} \quad A(x,y,0) = A_{i}(x,y)\exp{i \Phi(x,y)}
\end{split}
\end{equation}
and $A_{i}(x,y)$ is of finite support, \ie $A_{i}(x,y) = 0$ when $\sqrt{x^2 + y^2} > \frac{D_{i}}{l_0}$.


\subsubsection{Paraxial Approximation }

Letting $A = \exp{i k L_0 n_0 z} U(x,y,z)$ gives
\begin{equation}\label{equ:helmholtz_equation_3}
\begin{split}
  & \frac{l_0^2}{L_0^2} \left( \ppder{U}{z}  + 2i k L_0 n_0 \pder{U}{z} - k^2 L_0^2 n_0^2 U \right) 
  + \nabla_{\perp}^2 U + l_0^2 k^2 n_0^2 U + 2 l_0^2 k^2 n_0 \sigma_n n(x,y,z) U 
  + l_0^2 k^2 \sigma_n^2 n^2(x,y,z) U = 0 \\
  &\quad  \mbox{for}   
  \quad -\infty < x,y < \infty, \quad 0 < z < \frac{L}{L_0} 
  \quad \mbox{where} \quad U(x,y,0) = U_{i}(x,y) \exp{i\Phi(x,y)}
\end{split}
\end{equation}
and $U_{i}(x,y)$ is of finite support, \ie $U_{i}(x,y) = 0$ when 
$\sqrt{x^2 + y^2} > \frac{D_{i}}{l_0}$. 

\begin{equation}
\begin{split}
  & \frac{l_0^2}{L_0^2} \ppder{U}{z} + 2i n_0 \frac{l_0^2 k}{L_0} \pder{U}{z} + \nabla_{\perp}^2 U 
   +  2n_0  l_0^2 k^2 \sigma_n n(x,y,z)U + l_0^2 k^2 \sigma_n^2 n^2(x,y,z) U = 0 \\
  &\quad  \mbox{for}   
  \quad -\infty < x,y < \infty, \quad 0 < z < \frac{L}{L_0} \quad  \mbox{where} \quad U(x,y,0) = U_{i}(x,y)\exp{i \Phi(x,y)}
\end{split}
\end{equation}
and $U_{i}(x,y)$ is of finite support, \ie $U_{i}(x,y) = 0$ when $\sqrt{x^2 + y^2} > \frac{D_{i}}{l_0}$.

\begin{table}[htp]
\begin{center}
\begin{tabular}{|c|c|c|}
\hline
Ratio                      & Symbols                          & Range        \\       
\hline
\hline
Index Variance             & $\sigma_n =  C_0^{1/3} \frac{C_n L_0^{1/3}}{\sqrt{2}}$   & $ (10^{-9}\leftrightarrow 10^{-5})$ \\ 
\hline
Inner-scale/Outer-scale    & $\epsilon = \frac{l_0}{L_0}$     & $10^{-4} \leftrightarrow 10^{-6}\quad$    \\                  
\hline
Inner-scale/Fresnel Length &  $\delta = \frac{l_0^2 k }{L_0} = (1 \leftrightarrow 10^{2})(10^{-3}\leftrightarrow 10^{-1}) $ & $10^{-3} \leftrightarrow 10$  \\
\hline
Inner-scale/Wavelength     & $\gamma = l_0^2 k^2 \sigma_n$    & $10^{-3} \leftrightarrow 1$  \\
\hline
Inner-scale/Wavelength     & $\nu   = \gamma \sigma_n$        & $10^{-12} \leftrightarrow 10^{-5}$  \\ 
\hline   
Propagation Distance/Outer-Scale & $\Lambda = \frac{L}{L_0}$  & $1 \leftrightarrow 10$ \\
\hline
Aperture Diameter/Inner-Scale    & $\Delta_{\perp} = \frac{\tilde{D_{i}}}{l_0}$  & $10^3 \leftrightarrow 10$ \\
\hline
\end{tabular}
\end{center}
\caption{Dimensional Scales}
\label{Dimensional Scales}
\end{table}

\begin{equation}
 \epsilon^2 \ppder{U}{z} + 2i n_0 \delta \pder{U}{z}   + \nabla_{\perp}^2 U 
   +  2n_0  \gamma n(\vect{r}) U + \nu n^2(\vect{r}) U = 0 
\end{equation}
for $-\infty < x,y < \infty$, $0 < z < \Lambda$ where $U(x,y,0) = U_{i}(x,y)\exp{i \Phi(x,y)}$
and $U_{i}(x,y)$ is of finite support, i.e. $U_{i}(x,y) = 0$ when $\sqrt{x^2 + y^2} > \Delta_{\perp}$.
\begin{equation}
  2i n_0 \delta \pder{U}{z}  + \nabla_{\perp}^2 U + 2 n_0 \gamma n(\vect{r}) U = 0 
\end{equation}
for $-\infty < x,y < \infty$, $0 < z < \Lambda$ where $U(x,y,0) = U_{i}(x,y)\exp{i \Phi(x,y)}$
and $U_{i}(x,y)$ is of finite support, i.e. $U_{i}(x,y) = 0$ when $\sqrt{x^2 + y^2} > \Delta_{\perp}$.
Also, 
\begin{equation}
\begin{split}
  \cor{n_1,n_2} = 1 - \frac{1}{C_0^{2/3}} \left( (z_{1} - z_{2})^{2} 
  + \epsilon \left|\vect{r}_{\perp,1}-\vect{r}_{\perp,2}\right|^{2}\right)^{1/3}, 
  \quad  \sqrt{\epsilon} \leq \sqrt{(z_{1} - z_{2})^{2} + \epsilon \left|\vect{r}_{\perp,1}-\vect{r}_{\perp,2}\right|^{2}}  \leq  1
\end{split}
\end{equation}

\begin{equation}
  \cor{n_1,n_2} = 1 - 
    \frac{1}{C_0^{2/3}}  \begin{cases}
        \frac{1}{\epsilon^{2/3}} \left( (z_{1} - z_{2})^{2} 
  + \epsilon \left|\vect{r}_{\perp,1}-\vect{r}_{\perp,2}\right|^{2}\right) & 0 \leq \sqrt{(z_{1} - z_{2})^{2} + \epsilon \left|\vect{r}_{\perp,1}-\vect{r}_{\perp,2}\right|^{2}}  \leq  \sqrt{\epsilon} \\
        \left( (z_{1} - z_{2})^{2} 
  + \epsilon \left|\vect{r}_{\perp,1}-\vect{r}_{\perp,2}\right|^{2}\right)^{1/3} &  \sqrt{\epsilon} \leq \sqrt{(z_{1} - z_{2})^{2} + \epsilon \left|\vect{r}_{\perp,1}-\vect{r}_{\perp,2}\right|^{2}}  \leq  1  \\
   \end{cases}
\end{equation}



\subsubsection{Numerically}
If $z_{1} = z_{2}$
\begin{equation}
  \cor{n_1,n_2} = 1 - \frac{1}{C_0^{2/3}} \epsilon^{1/3} \left|\vect{r}_{1}-\vect{r}_{2}\right|^{2} \quad 0 \leq \left|\vect{r}_{1}-\vect{r}_{2}\right| \leq 1  
\end{equation}
\begin{equation}
  \cor{n_1,n_2} = 1 - \frac{1}{C_0^{2/3}} \epsilon^{1/3} \left|\vect{r}_{1}-\vect{r}_{2}\right|^{2/3} \quad 1 \leq \left|\vect{r}_{1}-\vect{r}_{2}\right| \leq \frac{1}{\sqrt{\epsilon}}  
\end{equation}

If $1 \geq |z_{1} - z_{2}| \geq \sqrt{\epsilon}$
\begin{equation}
    \cor{n_1,n_2} = 1 - \frac{1}{C_0^{2/3}} \left( (z_{1} - z_{2})^{2}  + \epsilon \left|\vect{r}_{1}-\vect{r}_{2}\right|^{2}\right)^{1/3} \\  
\end{equation}

\section{Lagrangian Approach}
The Lagrangian approach to atmospheric propagation involves recasting the propagation equation as a functional, where stationary points of the functional are also solutions of the propagation equation.  \href{https://en.wikipedia.org/wiki/Hamilton%27s_principle}{See Hamilton's Principle}

The functional for the paraxial equation above is defined as 
\begin{equation}
    J(u,\nabla u) = \int_{0}^{\Lambda} \iint \mathcal{L}_{D}(u,\nabla u) dx dy \, dz
\end{equation}
where $\mathcal{L}_{D}(u,\nabla u)$ is referred to as the Lagrangian Density, and is defined as 
\begin{equation}
\begin{split}
    \mathcal{L}_{D}(u,\nabla u) &= i n_0 \alpha  \left(\partial_{z}(u) \overline{u} - \overline{\partial_{z}(u)}u\right) - |\nabla_{\perp} u|^2 + 2n_0 \gamma^2 n(\vect{r}_{\perp},z) |u|^2 \\
    &= -2 n_0 \alpha\, \mathrm{Im}\left(\partial_{z}(u) \overline{u}\right) - |\nabla_{\perp} u|^2 + 2n_0 \gamma^2 n(\vect{r}_{\perp},z) |u|^2 
\end{split}
\end{equation}
It is fairly easy to show that critical points (points where derivative of the functional is zero) of this functional are also solutions of the propagation equation.  Derivatives of a functional are calculated much the same way an derivatives of a vector equation:
\begin{equation}
\begin{split}
    \lim_{\epsilon \to 0} &\frac{J(u+\epsilon v, \nabla u + \epsilon \nabla v)-J(u,\nabla u)}{\epsilon} \\
    &= \lim_{\epsilon \to 0} \frac{1}{\epsilon} \left[\int_{0}^{\Lambda} \iint \mathcal{L}_{D}(u+\epsilon v,\nabla u + \epsilon \nabla v) dx dy \, dz - \int_{0}^{\Lambda} \iint \mathcal{L}_{D}(u,\nabla u) dx dy \, dz\right] \\
    &= \int_{0}^{\Lambda} \iint \pder{\mathcal{L}_{D}}{u} v+ \pder{\mathcal{L}_{D}}{(\partial_{x}(u))} \partial_{x}(v) + \pder{\mathcal{L}_{D}}{(\partial_{y}(u))} \partial_{y}(v)   + \pder{\mathcal{L}_{D}}{(\partial_{z}(u))} \partial_{z}(v) \,dx dy \, dz\\
    &= \int_{0}^{\Lambda} \iint \left[\pder{\mathcal{L}_{D}}{u}  - \pder{}{x}\left(\pder{\mathcal{L}_{D}}{(\partial_{x}(u))}\right)
    -\pder{}{y}\left(\pder{\mathcal{L}_{D}}{(\partial_{y}(u))}\right)
    -\pder{}{z}\left(\pder{\mathcal{L}_{D}}{(\partial_{z}(u))}\right)\right] v \,dx dy \, dz \\
    &+ \pder{\mathcal{L}_{D}}{(\partial_{x}(u))}v \big|_{-\infty}^{\infty} + \pder{\mathcal{L}_{D}}{(\partial_{y}(u))} v \big|_{-\infty}^{\infty} + \pder{\mathcal{L}_{D}}{(\partial_{z}(u))} v \big|_{0}^{\Lambda}.
\end{split}
\end{equation}
We can neglect the boundary terms by assuming that $v(0) = v(\Lambda) = \lim_{x\to\pm\infty}v =\lim_{y\to\pm\infty}v =0$.  This leaves 
\begin{equation}
    \pder{J}{u} = \int_{0}^{\Lambda} \iint \left[\pder{\mathcal{L}_{D}}{u}  - \pder{}{x}\left(\pder{\mathcal{L}_{D}}{(\partial_{x}(u))}\right)
    -\pder{}{y}\left(\pder{\mathcal{L}_{D}}{(\partial_{y}(u))}\right)
    -\pder{}{z}\left(\pder{\mathcal{L}_{D}}{(\partial_{z}(u))}\right)\right] v \,dx dy \, dz 
\end{equation}
and thus stationary points are given by $u$ such that 
\begin{equation}
    \pder{J}{u} = \int_{0}^{\Lambda} \iint \left[\pder{\mathcal{L}_{D}}{u}  - \pder{}{x}\left(\pder{\mathcal{L}_{D}}{(\partial_{x}(u))}\right)
    -\pder{}{y}\left(\pder{\mathcal{L}_{D}}{(\partial_{y}(u))}\right)
    -\pder{}{z}\left(\pder{\mathcal{L}_{D}}{(\partial_{z}(u))}\right)\right] v \,dx dy \, dz = 0.
\end{equation}
However, the only way this holds for all $v$ is if 
\begin{equation}
    \pder{\mathcal{L}_{D}}{u}  - \pder{}{x}\left(\pder{\mathcal{L}_{D}}{(\partial_{x}(u))}\right)
    -\pder{}{y}\left(\pder{\mathcal{L}_{D}}{(\partial_{y}(u))}\right)
    -\pder{}{z}\left(\pder{\mathcal{L}_{D}}{(\partial_{z}(u))}\right) = 0.
\end{equation}
This is called the Euler-Lagrange equation(s) and one can verify that inserting the definition for $\mathcal{L}_{D}$ recovers the paraxial equation discussed above.  Thus, stationary points of $J$ are also solutions to the paraxial propagation equation.


\subsection{Variational Approximations for Low Dimensional Propagation}

One could proceed in a straight forward manor and try to find stationary points of the functional $J$ through numerical algorithms.  However, the approach we take here is to use the Lagrangian formulation to find low dimensional (finite dimensional) approximations to the governing equations through variational means.  The application of **variational approximations** to PDEs is not a new concept.  I was introduced to the methodology while studying the propagation of nonlinear solitary waves (solitons) in optical fiber.  The seminal paper for this application was [Anderson, *Variational approach to nonlinear pulse propagation in optical fibers*, (1983)], which introduced the method by following the work done in applying VA methods to soliton propagation in plasmas.  The method was further developed for the context of optical fiber in [Anderson, Lisak and Reichel, *Approximate analytical approaches to nonlinear pulse propagation in optical fibers: A comparison*, (1988)].  

The approximation procedure begins with postulation of an ansatz (trial solution in  analytically form) which has a given function for the solution which allows us to integrate out the $x$ and $y$ dimensions and use the corresponding Euler-Lagrange equations to derive stochastic ODEs for the $z$ evolution.


\subsection{Derivation of Gaussian Beam Solution}
Consider the unperturbed Helmholtz equation $(n(x,y,z) = 0)$
\begin{equation}
   2i n_0 \alpha \pder{U}{z}  + \nabla_{\perp}^2 U  = 2i n_0 \alpha \pder{U}{z}  + \frac{1}{r}\pder{}{r}\left(r\pder{U}{r}\right) + \frac{1}{r^2}\pder{U}{\theta^2} = 0 
\end{equation}
Now, lets look for a radially symmetric solution of the form 
\begin{equation}
\begin{split}
   U  \propto \exp{-i\left[ \frac{k r^2}{2 q(z)} + p(z) \right]}
\end{split}
\end{equation}
where $k=n_0 \alpha $.  The terms become
\begin{equation}
   \pder{U}{z} =  -i\left[-\frac{k r^2}{2 q^2(z)} \dot{q} + \dot{p}\right] U 
\end{equation}
\begin{equation}
   \frac{1}{r}\pder{}{r}\left(r\pder{U}{r}\right) =  -\left[i \frac{2k}{q} + \frac{k^2}{q^2} r^2\right] U 
\end{equation}
and the equation becomes 
\begin{equation}
   \begin{split}
       &2i k (-i) \left[-\frac{k r^2}{2 q^2} \dot{q} + \dot{p}\right] - \left[i \frac{2k}{q} + \frac{k^2}{q^2} r^2\right]   = \left[-\frac{2k^2 r^2}{2 q^2} \dot{q} + 2k\dot{p}\right] - \left[i \frac{2k}{q} + \frac{k^2}{q^2} r^2\right]  = \\
       &\left(-\frac{2k^2 r^2}{2 q^2} \dot{q} - i \frac{2k}{q}\right) + \left(2k\dot{p} -  \frac{k^2}{q^2} r^2\right) = -\frac{k^2 r^2}{q^2} \left(\dot{q} + 1\right) + 2k\left(\dot{p}- i \frac{1}{q}\right) = 0
   \end{split}
\end{equation}
Now, since this equation has to hold for all $r$, both terms above have to be independently equal to $0$.   
\begin{equation}
   \begin{split}
       \dot{q} + 1 &= 0 \rightarrow q = q(z_0) - (z-z_0)\\
       \dot{p}- i \frac{1}{q} &= 0 \rightarrow \dot{p} =  i \frac{1}{q(z_0) - (z-z_0)}
   \end{split}
\end{equation}
For the solution to be of finite engergy ($L_2$), the value of $q_0$ must be complex.  To see this consider if $q_0$ was real, then $|\exp{-i\left[ \frac{k r^2}{2 q(z)} + p(z) \right]}| = |\exp{-i p(z)}|$, which doesn't depend on $r$ and hence has infinite energy.  In addition, we can absorb any real part of $q(z_0)$ into $z_0$.  Hence, $q(z_0) = iq_0$ where $q_0$ is real.

For the $p$ equation, we have
\begin{equation}
   \begin{split}
       p(z) &= i \int_{z_0}^{z} \frac{1}{i q_0 - (z'-z_0)} dz'= -i \int_{0}^{z-z_0} \frac{1}{ z' - i q_0} dz' \\
       &= -i \ln{z' - i q_0}|_{0}^{z-z_0} = -i \ln{z - z_0 - i q_0} + i \ln{- i q_0}\\
       &= -i \left(\ln{z - z_0 - i q_0} - \ln{- i q_0}\right)\\
       &= -i \ln{\frac{z - z_0 - i q_0}{- i q_0}} = -i \ln{1 + i \frac{z - z_0}{q_0}}
   \end{split}
\end{equation}
so
\begin{equation}
   \begin{split}
       \exp{-i p(z)} = \exp{- \ln{1 + i \frac{z - z_0}{q_0}}}  = \frac{1}{1 + i \frac{z - z_0}{q_0}}
   \end{split}
\end{equation}

Now, we can represent complex numbers $z = a+ib$ as $(a^2+b^2)^{1/2}\exp{i \atan{\frac{b}{a}} }$ so 
\begin{equation}
   \begin{split}
       \exp{-i p(z)} = \frac{1}{\left(1 + \frac{(z - z_0)^2}{q_0^2}\right)^\frac{1}{2} \exp{i \atan{ \frac{z - z_0}{q_0} } }} = \left(1 + \frac{(z - z_0)^2}{q_0^2}\right)^{-\frac{1}{2}} \exp{-i \atan{ \frac{z - z_0}{q_0} } }
   \end{split} 
\end{equation}

Finally,
\begin{equation}
\begin{split}
   U  &\propto \left(1 + \frac{(z - z_0)^2}{q_0^2}\right)^{-\frac{1}{2}} \exp{i\left[ \frac{k r^2}{2 ((z'-z_0) - i q_0)} \right]} \exp{-i \atan{ \frac{z - z_0}{q_0} } } \\
   &\propto \left(1 + \frac{(z - z_0)^2}{q_0^2}\right)^{-\frac{1}{2}} \exp{i\left[ \frac{k r^2((z'-z_0) + i q_0)}{2 ((z'-z_0)^2 + q_0^2)} \right]} \exp{-i \atan{ \frac{z - z_0}{q_0} } } \\
   &\propto \left(1 + \frac{(z - z_0)^2}{q_0^2}\right)^{-\frac{1}{2}} \exp{i\left[ \frac{k r^2(z'-z_0)}{2 ((z'-z_0)^2 + q_0^2)} \right]} \exp{-\left[ \frac{k r^2 q_0}{2 ((z'-z_0)^2 + q_0^2)} \right]} \exp{-i \atan{ \frac{z - z_0}{q_0} } }\\
   &\propto \left(1 + \frac{(z - z_0)^2}{q_0^2}\right)^{-\frac{1}{2}} \exp{i\left[ \frac{k r^2}{2 (z'-z_0)\left(1 + \left(\frac{q_0}{z'-z_0}\right)^2\right)} \right]} \exp{-\left[ \frac{k r^2}{2q_0\left(1+\left(\frac{z'-z_0}{q_0}\right)^2\right)} \right]} \exp{-i \atan{ \frac{z - z_0}{q_0} } }\\
\end{split}
\end{equation}

Now we can define the spot size and beam waist as 
\begin{equation}
\begin{split}
   w(z) = w_0 \sqrt{1+\left(\frac{z-z_0}{q_0}\right)^2}, \quad  w_0 = \sqrt{q_0}
\end{split}
\end{equation}
so
\begin{equation}
\begin{split}
   w^2(z) = q_0\left(1+\left(\frac{z-z_0}{q_0}\right)^2\right) = \frac{1}{q_0} \left(q_0^2+(z-z_0)^2\right) \rightarrow q_0 w^2(z) = \left(q_0^2+(z-z_0)^2\right) 
\end{split}
\end{equation}


We can also define radius of curvature as 
\begin{equation}
\begin{split}
   R(z) = (z-z_0)\left[1 + \left(\frac{q_0}{ z-z_0}\right)^2\right]
\end{split}
\end{equation}
so
\begin{equation}
\begin{split}
  \frac{k r^2(z'-z_0)}{2 ((z'-z_0)^2 + q_0^2)} &= \frac{k r^2(z'-z_0)}{2 (z'-z_0)^2(1 + \frac{q_0^2}{(z'-z_0)^2})}\\
  &=\frac{k r^2}{2 R(z)}
\end{split}
\end{equation}

Finally,
\begin{equation}
\begin{split}
   U  &\propto \frac{w_0}{w(z)} \exp{i k \frac{r^2}{2 R(z)}} \exp{-k \frac{r^2}{2 w^2(z)}} \exp{-i \atan{ \frac{z - z_0}{q_0} } }
\end{split}
\end{equation}
Simplest form is to let $z_0 = 0$ , so the spot size is smallest at the origin, and put $k = n_0 \alpha$ and we get
\begin{equation}
\begin{split}
   U  = A_0\frac{w_0}{w(z)} \exp{in_0 \alpha\frac{r^2}{2 R(z)}} \exp{-n_0 \alpha\frac{r^2}{2w^2(z)}} \exp{-i \atan{ \frac{z}{q_0} } }
\end{split}
\end{equation}
where
\begin{equation}
\begin{split}
   R(z) = z\left[1 + \left(\frac{q_0}{z}\right)^2\right]\quad    w(z) = w_0 \sqrt{1+\left(\frac{z}{q_0}\right)^2} , \quad  w_0 = \sqrt{q_0}
\end{split}
\end{equation}
or 
\begin{equation}
\begin{split}
   q_0= w_0^2 \rightarrow R(z) = z\left[1 + \left(\frac{w_0^2}{z}\right)^2\right] 
\end{split}
\end{equation}
and
\begin{equation}
\begin{split}
   q_0= w_0^2 \rightarrow  w(z) = w_0 \sqrt{1+\left(\frac{z}{w_0^2}\right)^2}
\end{split}
\end{equation}


Alternatively, for our purposes we need $z_0 = \Lambda$
\begin{equation}
\begin{split}
   R(z) = (z-\Lambda)\left[1 + \left(\frac{w_0^2}{(z-\Lambda)}\right)^2\right]
\end{split}
\end{equation}
and 
\begin{equation}
\begin{split}
   R(0) = -\Lambda\left[1 + \left(\frac{w_0^2}{\Lambda}\right)^2\right]
\end{split}
\end{equation}

\begin{equation}
\begin{split}
   w(z) = w_0 \sqrt{1+\left(\frac{(z-\Lambda)}{w_0^2}\right)^2}
\end{split}
\end{equation}
and
\begin{equation}
\begin{split}
   w(0) = w_0 \sqrt{1+\left(\frac{\Lambda}{w_0^2}\right)^2}
\end{split}
\end{equation}

So in total...
\begin{equation}
\begin{split}
   U  = A_0\frac{w_0}{w(z)} \exp{in_0 \alpha\frac{r^2}{2 R(z)}} \exp{-n_0 \alpha\frac{r^2}{2w^2(z)}} \exp{-i \atan{ \frac{z}{w_0^2} } }
\end{split}
\end{equation}
\begin{equation}
\begin{split}
   w(z) = w_0 \sqrt{1+\left(\frac{z-z_0}{w_0^2}\right)^2} 
   = \frac{w_0}{w_0^2} \sqrt{w_0^4+\left(z-z_0\right)^2} = \frac{1}{w_0} \sqrt{w_0^4+\left(z-z_0\right)^2}
\end{split}
\end{equation}
\begin{equation}
\begin{split}
   R(z) = (z-z_0)\left[1 + \left(\frac{w_0^2}{z-z_0}\right)^2\right] = \frac{1}{z-z_0} \left[(z-z_0)^2 + w_0^4\right] 
\end{split}
\end{equation}
But I would rather work with the inverse of these
\begin{equation}
\begin{split}
   W(z) = \frac{1}{w(z)} = \frac{w_0}{\sqrt{w_0^4+\left(z-z_0\right)^2}} 
\end{split}
\end{equation}
\begin{equation}
\begin{split}
   F(z) = \frac{1}{R(z)} = \frac{z-z_0}{(z-z_0)^2 + w_0^4}
\end{split}
\end{equation}
with solution 
\begin{equation}
\begin{split}
   U  = A_0 w_0 W(z) \exp{i\frac{n_0 \alpha}{2} F(z) r^2} \exp{- \frac{n_0 \alpha}{2} W^2(z) r^2} \exp{-i \atan{ \frac{z}{w_0^2} } }
\end{split}
\end{equation}

%\subsubsection{Putting Dimension in Back In}
%
%Letting $z = \dim{z}/L_0, r = \dim{r}/l_0$ and $q_0 =\dim{q_0}/l_0$ recall 
%$\alpha = \frac{l_0^2 k }{L_0}$
%
%\begin{equation}
%\begin{split}
%   U  = A_0\frac{w_0}{w(\dim{z})} \exp{in_0\alpha \frac{1}{l_0^2} \frac{\dim{r}^2}{2 R(\dim{z})}} \exp{-n_0\alpha \frac{1}{l_0^2}\frac{\dim{r}^2}{2 w^2(\dim{z})}} \exp{-i \atan{ \frac{\dim{z}}{L_0 q_0} } }
%\end{split}
%\end{equation}
%where
%\begin{equation}
%\begin{split}
%   R(\dim{z}) &= \frac{\dim{R}}{l_0} = (\dim{z}/L_0-\Lambda)\left[1 + \left(\frac{q_0}{(\dim{z}/L_0-\Lambda)}\right)^2\right] \\
%   &= \frac{1}{L_0} (\dim{z}-L)\left[1 + \left(\frac{L_0 w_0^2}{(\dim{z}-L)}\right)^2\right] \\
%   &= \frac{1}{L_0} (\dim{z}-L)\left[1 + \left(\frac{L_0 q_0}{(\dim{z}-L)}\right)^2\right] \\
%\end{split}
%\end{equation}
%\begin{equation}
%\begin{split}
%   R(\dim{z}) &= (\dim{z}/L_0-\Lambda)\left[1 + \left(\frac{q_0}{(\dim{z}/L_0-\Lambda)}\right)^2\right] \\
%   &= \frac{1}{L_0} (\dim{z}-L)\left[1 + \left(\frac{L_0 w_0^2}{(\dim{z}-L)}\right)^2\right] \\
%   &= \frac{1}{L_0} (\dim{z}-L)\left[1 + \left(\frac{L_0 q_0}{(\dim{z}-L)}\right)^2\right] \\
%   &= \frac{1}{L_0} \dim{R}(\dim{z})
%\end{split}
%\end{equation}
%and 
%\begin{equation}
%\begin{split}
%   R(0) = -L\left[1 + \left(\frac{L_0 q_0}{L}\right)^2\right]
%\end{split}
%\end{equation}
%
%\begin{equation}
%\begin{split}
%   w(\dim{z}) &= \sqrt{q_0} \sqrt{1+\left(\frac{(\dim{z}-L)}{L_0 q_0}\right)^2}\\
%              &= \sqrt{q_0}\frac{1}{L_0 q_0} \sqrt{(L_0 q_0)^2+(\dim{z}-L)^2}\\
%              &= \frac{1}{L_0} \sqrt{(L_0 q_0)^2+(\dim{z}-L)^2}\\
%              &= \frac{1}{L_0} \dim{w}(\dim{z})
%\end{split}
%\end{equation}
%and
%\begin{equation}
%\begin{split}
%   w(0) = \frac{1}{L_0} \dim{w}(0) = \frac{1}{L_0} \sqrt{(L_0 q_0)^2+ L^2}
%\end{split}
%\end{equation}
%
%
%\begin{equation}
%\begin{split}
%   U  &= A_0\frac{L_0w_0}{\dim{w}(\dim{z})} \exp{in_0\alpha \frac{L_0}{l_0^2} \frac{\dim{r}^2}{2 \dim{R}(\dim{z})}} \exp{-n_0\alpha \frac{L_0^2}{l_0^2}\frac{\dim{r}^2}{2 \dim{w}^2(\dim{z})}} \exp{-i \atan{ \frac{\dim{z}}{L_0 q_0} } }\\
%   &=A_0\frac{L_0w_0}{\dim{w}(\dim{z})} \exp{in_0\frac{l_0^2 k }{L_0} \frac{L_0}{l_0^2} \frac{\dim{r}^2}{2 \dim{R}(\dim{z})}} \exp{-n_0\frac{l_0^2 k }{L_0} \frac{L_0^2}{l_0^2}\frac{\dim{r}^2}{2 \dim{w}^2(\dim{z})}} \exp{-i \atan{ \frac{\dim{z}}{L_0 q_0} } }\\
%   &=A_0\frac{L_0w_0}{\dim{w}(\dim{z})} \exp{in_0 k  \frac{\dim{r}^2}{2 \dim{R}(\dim{z})}} \exp{-n_0 k L_0\frac{\dim{r}^2}{2 \dim{w}^2(\dim{z})}} \exp{-i \atan{ \frac{\dim{z}}{L_0 q_0} } }\\  
%\end{split}
%\end{equation}
%
\subsubsection{Normalization}
\begin{equation}
   |U|^2  = A^2_0 \frac{w_0^2}{w^2(z)} \exp{- n_0 \alpha \frac{r^2}{w^2(z)}} 
\end{equation}
so
\begin{equation}
\iint |U|^2 \, dx \, dy = 2 \pi A^2_0 \frac{w_0^2}{w^2(z)}  \int_{0}^{\infty}\exp{- n_0 \alpha\frac{r^2}{w^2(z)}} r \,dr
\end{equation}
letting $r = \frac{w(z)}{\sqrt{ n_0 \alpha }} x$ so $dr = \frac{w(z)}{\sqrt{ n_0 \alpha }} dx$
\begin{equation}
2 \pi A^2_0 \frac{w_0^2}{w^2(z)}  \int_{0}^{\infty}\exp{-r^2} \frac{w(z)}{\sqrt{ n_0 \alpha }} x \,\frac{w(z)}{\sqrt{ n_0 \alpha }} dx
\end{equation}

\begin{equation}
2 \pi A^2_0 \frac{w_0^2}{n_0 \alpha }  \int_{0}^{\infty}\exp{-x^2}  x \, dx = \pi A^2_0 \frac{w_0^2}{n_0 \alpha } = 1
\end{equation}

\begin{equation}
A_0 = \sqrt{\frac{n_0 \alpha }{\pi w_0^2}} = \sqrt{\frac{n_0 \alpha }{\pi}} \frac{1}{w_0}
\end{equation}


\begin{equation}
\begin{split}
   W(z) = \frac{1}{w(z)} = \frac{w_0}{\sqrt{w_0^4+\left(z-z_0\right)^2}} 
\end{split}
\end{equation}
\begin{equation}
\begin{split}
   F(z) = \frac{1}{R(z)} = \frac{z-z_0}{(z-z_0)^2 + w_0^4}
\end{split}
\end{equation}
with solution 
\begin{equation}
\begin{split}
   U  = \sqrt{\frac{n_0 \alpha}{\pi}} W(z) \exp{i\frac{n_0 \alpha}{2} F(z) r^2} \exp{- \frac{n_0 \alpha}{2} W^2(z) r^2} \exp{-i \atan{ \frac{z}{w_0^2} } }
\end{split}
\end{equation}











%\section{Perturbed Gaussian Beam I}
We begin by using a Gaussian ansatz:
\begin{equation}
 u(\vect{r}_{\perp},A,W,R,P) = \frac{A}{W} \exp{-n_0 \alpha \frac{r^2}{2W^2}} \exp{ i n_0 \alpha \frac{r^2}{2R}} \exp{- i P }
\end{equation}
where $r = |\vect{r}_{\perp}|$ which is motivated by the exact Gaussian beam solution discussed above.   Each of the parameters $A,W,R,P$ are functions of $z$.  Now, lets calculate each term in the Lagrangian density:
\begin{equation}
    \begin{split}
    \partial_{z}(u) &= \left[ \frac{\dot{A}}{A} - \frac{\dot{W}}{W} + n_0 \alpha \frac{r^2}{W^3} \dot{W} - i n_0 \alpha \frac{r^2}{2 R^2} \dot{R} - i \dot{P}  \right] u\\
     \partial_{z}(u)\conj{u} &= \left[ \frac{\dot{A}}{A} - \frac{\dot{W}}{W} + n_0 \alpha \frac{r^2}{W^3} \dot{W} - i n_0 \alpha \frac{r^2}{2 R^2} \dot{R} - i \dot{P}  \right] |u|^2 \\
     \partial_{z}(u)\conj{u} &=  \left[ \frac{\dot{A}}{A} - \frac{\dot{W}}{W} + n_0 \alpha \frac{r^2}{W^3} \dot{W} - i n_0 \alpha \frac{r^2}{2 R^2} \dot{R} - i \dot{P}  \right] \frac{A^2}{W^2}\exp{- \frac{r^2}{W^2}}\\
     \overline{\partial_{z}(u)}u &=  \left[ \frac{\dot{A}}{A} - \frac{\dot{W}}{W} + n_0 \alpha \frac{r^2}{W^3} \dot{W} + i n_0 \alpha \frac{r^2}{2 R^2} \dot{R} + i \dot{P}  \right] \frac{A^2}{W^2}\exp{- \frac{r^2}{W^2}}\\     
     i n_0 \alpha  \left(\partial_{z}(u) \overline{u} - \overline{\partial_{z}(u)}u\right) &=
      i n_0 \alpha \left[ - i n_0 \alpha \frac{r^2}{ R^2} \dot{R} - i 2 n_0 \alpha \dot{P}  \right] \frac{A^2}{W^2}\exp{- \frac{r^2}{W^2}}\\ 
      &= \left[ (n_0 \alpha)^2 \frac{r^2}{ R^2} \dot{R} + 2 n_0 \alpha \dot{P}  \right] \frac{A^2}{W^2}\exp{- \frac{r^2}{W^2}}\\
    \end{split}
\end{equation}


\begin{equation}
    \begin{split}
    u_{x} &= n_0 \alpha \left(i \frac{1}{R} - \frac{1}{W^2}\right) x u\\
    |u_{x}|^2 &= (n_0 \alpha)^2 \left(\frac{1}{R^2} + \frac{1}{W^4}\right) x^2 |u|^2\\
    u_{y} &= n_0 \alpha \left(i \frac{1}{R} - \frac{1}{W^2}\right) y u\\
    |u_{y}|^2 &= (n_0 \alpha)^2\left(\frac{1}{R^2} + \frac{1}{W^4}\right) y^2 |u|^2\\
    |u_{x}|^2 +|u_{y}|^2 &= (n_0 \alpha)^2 \left(\frac{1}{R^2} + \frac{1}{W^4}\right) r^2 |u|^2\\ 
    \end{split}
\end{equation}


\begin{equation}
    \begin{split}
    |u|^2 &= \frac{A^2}{W^2} \exp{-r^2 W^2}   
    \end{split}
\end{equation}


Putting this into the definition of the Lagrangian Density gives
\begin{equation}
    \begin{split}
    \mathcal{L}_{D}(u,\nabla u) &=   \left[ (n_0 \alpha)^2 \frac{r^2}{ R^2} \dot{R} + 2 n_0 \alpha \dot{P}    - (n_0 \alpha)^2 \left(\frac{1}{R^2} + \frac{1}{W^4}\right) r^2 + 2n_0 \gamma^2 n(\vect{r}_{\perp},z) \right] |u|^2 \\   
    \end{split}
\end{equation}

and then the action becomes
\begin{equation}
    \begin{split}
        J(u,\nabla u) &= \int_{0}^{\Lambda} \iint \mathcal{L}_{D}(u,\nabla u) dx dy \, dz\\
        J(A,W,R,P;\dot{A},\dot{W},\dot{R},\dot{P}) &= \int_{0}^{\Lambda} (n_0 \alpha)^2 \frac{1}{ R^2} \dot{R} I_2 + 2 n_0 \alpha \dot{P} I_0 - (n_0 \alpha)^2 \left(\frac{1}{R^2} + \frac{1}{W^4}\right) I_2 + 2n_0 \gamma^2 \left<n(\vect{r}_{\perp},z), |u|^2\right>  dz\\ 
    \end{split}
\end{equation}
where is $\left<n(\vect{r}_{\perp},z), |u|^2\right>$ the projection of $|u|^2$ onto the index realizations $n(\vect{r}_{\perp},z)$ and 
\begin{equation}
    \begin{split}
        I_0 &=  \int_{0}^{2\pi} \int_{0}^{\infty} \frac{A^2}{W^2} \exp{-n_0 \alpha \frac{r^2 }{W^2}} r \,dr \,d\theta \\
        &=  2\pi \frac{A^2}{W^2} \int_{0}^{\infty}  \exp{-n_0 \alpha \frac{r^2}{W^2}} r \,dr \\
        &=  \frac{2\pi}{n_0 \alpha} A^2 \int_{0}^{\infty}  \exp{-x^2} x \,dx \\
        &= -\frac{\pi}{n_0 \alpha} A^2 \exp{-x^2}\big|_{0}^{\infty} \\
        &= \frac{\pi}{n_0 \alpha} A^2 
    \end{split}
\end{equation}
\begin{equation}
    \begin{split}
      I_2 &=  \int_{0}^{2\pi} \int_{0}^{\infty} r^2 \frac{A^2}{W^2} \exp{-n_0 \alpha \frac{r^2 }{W^2}} r \,dr \,d\theta\\           
          &= 2\pi\frac{A^2}{W^2} \int_{0}^{\infty} r^3 \exp{-n_0 \alpha \frac{r^2 }{W^2}} \,dr \\
          &= \frac{2\pi}{(n_0 \alpha)^2} W^2 A^2 \int_{0}^{\infty} x^3 \exp{-x^2} \,dx \\
          &= \frac{\pi}{(n_0 \alpha)^2} W^2 A^2 \int_{0}^{\infty} y \exp{-y} \,dy  \\
          &= \frac{\pi}{(n_0 \alpha)^2} W^2 A^2
    \end{split}
\end{equation}
or
\begin{equation}
    \begin{split}
        J(A,W,R,P;\dot{A},\dot{W},\dot{R},\dot{P}) &= \int_{0}^{\Lambda} (n_0 \alpha)^2 \frac{1}{ R^2} \dot{R} \frac{\pi}{(n_0 \alpha)^2} W^2 A^2 + 2 n_0 \alpha \dot{P} \frac{\pi}{n_0 \alpha} A^2 - (n_0 \alpha)^2 \left(\frac{1}{R^2} + \frac{1}{W^4}\right) \frac{\pi}{(n_0 \alpha)^2} W^2 A^2 + 2n_0 \gamma^2 \left<n(\vect{r}_{\perp},z), |u|^2\right>  dz\\ 
    &=\int_{0}^{\Lambda} F_{D}(A,W,R,P;\dot{A},\dot{W},\dot{R},\dot{P})  dz\\
    \end{split}
\end{equation}
where
\begin{equation}
    F_{D}(A,W,R,P;\dot{A},\dot{W},\dot{R},\dot{P}) = \pi \frac{W^2 A^2}{R^2} \dot{R} + 2 \pi A^2 \dot{P}   - \pi  A^2  \left(\frac{W^2}{R^2} + \frac{1}{W^2}\right) + 2n_0 \gamma^2 \left<n(\vect{r}_{\perp},z), |u|^2\right>  
\end{equation}
 
\begin{equation}
    F_{D}(A,W,R,P;\dot{A},\dot{W},\dot{R},\dot{P}) = A^2\pi \left[\frac{W^2}{R^2} \dot{R} + 2 \dot{P} - \left(\frac{W^2}{R^2} + \frac{1}{W^2}\right)\right] + 2n_0 \gamma^2 \left<n(\vect{r}_{\perp},z), |u|^2\right>  
\end{equation} 

\subsection{Euler-Lagrange equation(s)}

Now, we can apply the Euler-Lagrange equations to the reduced density $F_{D}(A,W,R,P;\dot{A},\dot{W},\dot{R},\dot{P})$, i.e.
\begin{equation}
\begin{split}
    \der{F_{D}}{A}  - \der{}{z}\left(\pder{F_{D}}{\dot{A}}\right) &= 0 \\
    \der{F_{D}}{P}  - \der{}{z}\left(\pder{F_{D}}{\dot{P}}\right) &= 0 \\
    \der{F_{D}}{R}  - \der{}{z}\left(\pder{F_{D}}{\dot{R}}\right) &= 0 \\
    \der{F_{D}}{W}  - \der{}{z}\left(\pder{F_{D}}{\dot{W}}\right) &= 0. \\
\end{split}
\end{equation}
So lets calculate the terms on the left side of each:
\begin{equation}
\begin{split}
    \der{F_{D}}{A}  &= 2A \pi \left[\left(\frac{W^2}{R^2}\dot{R}  + 2\dot{P}\right) -  \left(\frac{W^2}{R^2} + \frac{1}{W^2}\right) \right] + 2 n_0 \gamma^2 \left< n(\vect{r}_{\perp},z),\pder{|u|^2}{A}\right>\\
\der{}{z}\left(\pder{F_{D}}{\dot{A}}\right) &= \der{}{z}(0) = 0 \\
\end{split}
\end{equation}


\begin{equation}
\begin{split}
   \der{F_{D}}{P}  &= 2n_0 \gamma^2 \left< n(\vect{r}_{\perp},z),\pder{|u|^2}{P}  \right> = 0 \\
\der{}{z}\left(\pder{F_{D}}{\dot{P}}\right) &= \der{}{z}\left( 2 \pi A^2 \right)  \\
\end{split}
\end{equation}


\begin{equation}
\begin{split}
    \der{F_{D}}{W}  &=  2 \pi A^2 \left[\frac{W}{R^2} \dot{R} -  \frac{1}{R^2}W + \frac{1}{W^3}\right] +  2n_0 \gamma^2 \left< n(\vect{r}_{\perp},z),\pder{|u|^2}{W}  \right> \\
\der{}{z}\left(\pder{F_{D}}{\dot{W}}\right) &= \der{}{z}(0) = 0 \\
\end{split}
\end{equation}


\begin{equation}
\begin{split}
    \der{F_{D}}{R}  &=  A^2 \pi \left[-\frac{W^2}{R^3}\dot{R} + \frac{W^2}{R^3}\right] +  2n_0 \gamma^2 \left< n(\vect{r}_{\perp},z),\pder{|u|^2}{R}  \right> \\
\der{}{z}\left(\pder{F_{D}}{\dot{R}}\right) &= \der{}{z}\left(  \pi A^2 \frac{W^2}{R^2}\right)\\
&= \pi A^2 \left(\frac{2W }{R^2} \dot{W} - 2\frac{W^2}{R^3} \dot{R}\right)\\
&= 2\pi A^2 \left(\frac{W}{R^2} \dot{W} - \frac{W^2}{R^3} \dot{R}\right)\\
\end{split}
\end{equation}


\subsubsection{Equations in P}
Since $|u|^2$ is only dependent on $A,W$, not $R,P$, the last of these equations says
\begin{equation}
\begin{split}
    \der{}{z}\left( 2 \pi A^2 \right) = 0,
\end{split}
\end{equation}
or that
\begin{equation}
\begin{split}
    A(z) = A(0)
\end{split}
\end{equation}

\subsubsection{Equations in W}

\begin{equation}
\begin{split}
    \der{}{z}\left(\pder{F_{D}}{\dot{W}}\right) &= \der{F_{D}}{W} \\
    2 \pi A^2 \left[\frac{W}{R^2} \dot{R} -  \frac{1}{R^2}W + \frac{1}{W^3}\right] +  2n_0 \gamma^2 \left< n(\vect{r}_{\perp},z),\pder{|u|^2}{W}  \right>  &= 0\\
\end{split}
\end{equation}
\begin{equation}
\begin{split}
     \left[\frac{W}{R^2} \dot{R} -  \frac{1}{R^2}W + \frac{1}{W^3}\right] = -  \frac{2n_0 \gamma^2}{2 \pi A^2} \left< n(\vect{r}_{\perp},z),\pder{|u|^2}{W}  \right>  \\
     \frac{W}{R^2} \dot{R} =  \frac{1}{R^2}W - \frac{1}{W^3}  -  \frac{2n_0 \gamma^2}{2 \pi A^2} \left< n(\vect{r}_{\perp},z),\pder{|u|^2}{W}  \right>  \\
     \dot{R} =  \frac{R^2}{W} \left(\frac{1}{R^2}W + \frac{1}{W^3}\right)  -  \frac{R^2}{W}\frac{2n_0 \gamma^2}{2 \pi A^2} \left< n(\vect{r}_{\perp},z),\pder{|u|^2}{W}  \right>  \\
       \dot{R} =   \left(1 -  \frac{R^2}{W^4 }\right) -  \frac{n_0 \gamma^2}{\pi} \frac{R^2}{W A_0^2 } \left< n(\vect{r}_{\perp},z),\pder{|u|^2}{W}  \right> \\
\end{split}
\end{equation}

\subsubsection{Equations in R}
\begin{equation}
\begin{split}
   \der{}{z}\left( \pder{F_{D}}{\dot{R}} \right) &= \der{F_{D}}{R} \\
\der{}{z}\left( \pi \frac{W^2 A^2}{R^2}\right) &= - 2 \pi \frac{A^2 W^2}{R^3}\dot{R} + 2\pi \frac{A^2 W^2}{R^3}
\end{split}
\end{equation}
and because $\dot{A} = 0$ we get
\begin{equation}
\begin{split}
2 \frac{W}{R^2} \dot{W}  - 2 \frac{W^2}{R^3} \dot{R} &= -  2\frac{W^2}{R^3} \dot{R} + 2\frac{W^2}{R^3} \\
\frac{W}{R^2} \dot{W}  &=  \frac{W^2}{R^3} \\
\dot{W}  &=  \frac{R^2}{W} \frac{W^2}{R^3} \\
\dot{W}  &= \frac{W}{R} \\
\end{split}
\end{equation}

\subsubsection{Equations in A}
\begin{equation}
\begin{split}
    \der{F_{D}}{A}  - \der{}{z}\left(\pder{F_{D}}{\dot{A}}\right) &= 0 \\
   2A_0 \pi\left[ \left(\frac{W^2}{R^2}\dot{R} + 2 \dot{P}\right)  - \left(\frac{1}{W^2} + \frac{W^2}{R^2}\right) \right] +  2n_0 \gamma^2 \left< n(\vect{r}_{\perp},z),\pder{|u|^2}{A}\right>  &= 0 \\
   \end{split}
\end{equation}

\begin{equation}
\begin{split}
\left(\frac{W^2}{R^2}\dot{R} + 2 \dot{P}\right)  - \left(\frac{1}{W^2} + \frac{W^2}{R^2}\right) &= -  \frac{n_0 \gamma^2}{ A_0 \pi} \left< n(\vect{r}_{\perp},z),\pder{|u|^2}{A}\right>  \\
 2 \dot{P}  &= - \frac{W^2}{R^2}\dot{R} + \left(\frac{1}{W^2} + \frac{W^2}{R^2}\right) -  \frac{n_0 \gamma^2}{ A_0 \pi} \left< n(\vect{r}_{\perp},z),\pder{|u|^2}{A}\right>  \\
  \dot{P}  &= - \frac{W^2}{2R^2}\dot{R} + \frac{1}{2}\left(\frac{1}{W^2} + \frac{W^2}{R^2}\right) -  \frac{n_0 \gamma^2}{ 2 A_0 \pi} \left< n(\vect{r}_{\perp},z),\pder{|u|^2}{A}\right>  \\
   \end{split}
\end{equation}

Putting in the equation for $\dot{R}$
\begin{equation}
\begin{split}
  \dot{P}  &= - \frac{W^2}{2R^2}\left[\left(1 -  \frac{R^2}{W^4 }\right) -  \frac{n_0 \gamma^2}{\pi} \frac{R^2}{W A_0^2 } \left< n(\vect{r}_{\perp},z),\pder{|u|^2}{W}  \right> \right] + \frac{1}{2}\left(\frac{1}{W^2} + \frac{W^2}{R^2}\right) -  \frac{n_0 \gamma^2}{ 2 A_0 \pi} \left< n(\vect{r}_{\perp},z),\pder{|u|^2}{A}\right>  \\    
  \dot{P}  &= \left[\left(- \frac{W^2}{2R^2} +  \frac{W^2}{2R^2} \frac{R^2}{W^4 }\right) + \frac{W^2}{2R^2}  \frac{n_0 \gamma^2}{\pi} \frac{R^2}{W A_0^2 } \left< n(\vect{r}_{\perp},z),\pder{|u|^2}{W}  \right> \right] + \frac{1}{2}\left(\frac{1}{W^2} + \frac{W^2}{R^2}\right) -  \frac{n_0 \gamma^2}{ 2 A_0 \pi} \left< n(\vect{r}_{\perp},z),\pder{|u|^2}{A}\right>  \\   
 \dot{P}  &= - \frac{W^2}{2R^2} + \frac{1}{2W^2 } +  \frac{n_0 \gamma^2}{2\pi} \frac{W}{ A_0^2 } \left< n(\vect{r}_{\perp},z),\pder{|u|^2}{W}  \right>  + \frac{1}{2W^2} + \frac{W^2}{2R^2} -  \frac{n_0 \gamma^2}{ 2 A_0 \pi} \left< n(\vect{r}_{\perp},z),\pder{|u|^2}{A}\right>  \\ 
 \dot{P}  &=  \frac{1}{W^2 } +  \frac{n_0 \gamma^2}{2\pi} \frac{W}{ A_0^2 } \left< n(\vect{r}_{\perp},z),\pder{|u|^2}{W}  \right>   -  \frac{n_0 \gamma^2}{ 2 A_0 \pi} \left< n(\vect{r}_{\perp},z),\pder{|u|^2}{A}\right>  \\ 
  \dot{P}  &=  \frac{1}{W^2 } +  \frac{n_0 \gamma^2}{2\pi A_0} \left< n(\vect{r}_{\perp},z),\frac{W}{ A_0 } \pder{|u|^2}{W}  - \pder{|u|^2}{A}\right>  \\
\end{split}
\end{equation}


\subsubsection{Finding the modes that project onto the index realizations}


\begin{equation}
\begin{split}
  \pder{|u|^2}{W} &= \pder{}{W}\left(\frac{A^2}{W^2}\exp{-\frac{r^2}{W^2}}\right)\\
  &= -2\frac{A^2}{W^3} \exp{-\frac{r^2}{W^2}} + 2\frac{A^2}{W^2}\frac{r^2}{W^3}\exp{-\frac{r^2}{W^2}}\\
  &= 2\frac{1}{W}\left[\frac{r^2 }{W^2} - 1\right] \frac{A^2}{W^2} \exp{-\frac{r^2}{W^2}}\\
  &= 2\frac{1}{W}\left[\frac{r^2 }{W^2} - 1\right] |u|^2
\end{split}
\end{equation}

\begin{equation}
\begin{split}
  \pder{|u|^2}{A} &= \pder{}{A}\left(\frac{A^2}{W^2}\exp{-\frac{r^2}{W^2}}\right)\\
  &= 2\frac{A}{W^2} \exp{-\frac{r^2}{W^2}} \\
  &= \frac{2}{A} |u|^2
\end{split}
\end{equation}

\begin{equation}
\begin{split}
  \frac{W}{ A } \pder{|u|^2}{W}  - \pder{|u|^2}{A} &= 2\frac{W}{ A }  \frac{1}{W}\left[\frac{r^2 }{W^2} - 1\right] |u|^2 - \frac{2}{A} |u|^2\\
  &= 2\frac{1}{ A }  \left[\frac{r^2 }{W^2}|u|^2 - |u|^2 - |u|^2 \right]\\
  &= \frac{2}{A} \left[\frac{r^2 }{W^2} - 2  \right]|u|^2
\end{split}
\end{equation}

\begin{equation}
\begin{split}
      V_{R} &= \frac{n_0 \gamma^2 }{\pi} \frac{R^2}{W A_0^2 } \left< n(\vect{r}_{\perp},z),\pder{|u|^2}{W}  \right> \\
      &= \frac{n_0 \gamma^2 }{\pi} \frac{R^2}{W^2 A_0^2 } \left< n(\vect{r}_{\perp},z),\left[\frac{r^2 }{W^2} - 1\right] |u|^2 \right>
\end{split}
\end{equation}


 
 
\begin{equation}
\begin{split}
      V_{P} &=  \frac{n_0 \gamma^2}{2\pi A_0} \left< n(\vect{r}_{\perp},z),\frac{W}{ A_0 } \pder{|u|^2}{W}  - \pder{|u|^2}{A}\right> \\
      &=\frac{n_0 \gamma^2}{2\pi A_0} \left< n(\vect{r}_{\perp},z),\frac{2}{A} \left[\frac{r^2 }{W^2} - 2  \right]|u|^2\right> \\
      &=\frac{n_0 \gamma^2}{\pi A_0^2} \left< n(\vect{r}_{\perp},z), \left[\frac{r^2 }{W^2} - 2  \right]|u|^2\right> \\
\end{split}
\end{equation}



%\section{Gaussian Ansatz}

We begin by using a Gaussian ansatz:
\begin{equation}
    u(\vect{r}_{\perp},A,W,R,P) = \frac{A W}{\sqrt{\pi}} \exp{-\frac{r^2}{2}W^2 + i \frac{r^2}{2}R - i P }
\end{equation}
where $r = |\vect{r}_{\perp}|$ which is motivated by the exact Gaussian beam solution to the \href{https://en.wikipedia.org/wiki/Gaussian_beam}{paraxial wave equation}.   As described above, each of the parameters $A,W,R,P$ are functions of $z$.  Now, lets calculate each term in the Lagrangian density:
\begin{equation}
    \begin{split}
    \partial_{z}(u) &= \left[ \frac{\dot{A}}{A} + \frac{\dot{W}}{W} - r^2 W \dot{W} + i \frac{r^2}{2}\dot{R} - i \dot{P}  \right] u\\
    \partial_{z}(u)\conj{u} &= \left[ \frac{\dot{A}}{A} + \frac{\dot{W}}{W} - r^2 W \dot{W} + i \frac{r^2}{2}\dot{R} - i \dot{P}  \right] |u|^2 \\
    &= \left[ \frac{\dot{A}}{A} + \frac{\dot{W}}{W} - r^2 W \dot{W} + i \frac{r^2}{2}\dot{R} - i \dot{P}  \right] \frac{A^2W^2}{\pi} \exp{-r^2 W^2}\\
    -2 n_0 \alpha\, \mathrm{Im}\left(\partial_{z}(u) \conj{u}\right) &= n_0 \alpha\,\left[2 |u|^2 \dot{P}-r^2 |u|^2 \dot{R} \right]
    \end{split}
\end{equation}
\begin{equation}
    \begin{split}
    |\nabla_{\perp} u|^2 &= r^2 |u|^2\left(W^4 + R^2\right)\\
    -|\nabla_{\perp} u|^2 &= -r^2 |u|^2\left(W^4 + R^2\right)
    \end{split}
\end{equation}
\begin{equation}
    \begin{split}
    |u|^2 &= \frac{A^2W^2}{\pi} \exp{-r^2 W^2}   
    \end{split}
\end{equation}

\begin{equation}
    \begin{split}
    \int_{0}^{2\pi} \int_{0}^{\infty} |u|^2 r dr d\theta &= 2 \pi \frac{A^2W^2}{\pi} \int_{0}^{\infty} \exp{-r^2 W^2}  r dr   = 2 A^2 \int_{0}^{\infty} \exp{-\rho^2}  \rho d\rho =  A^2 
    \end{split}
\end{equation}

Putting this into the definition of the Lagrangian Density gives
\begin{equation}
    \begin{split}
    \mathcal{L}_{D}(u,\nabla u) &=  -2 n_0 \alpha\, \mathrm{Im}\left(\partial_{z}(u) \conj{u}\right) - |\nabla_{\perp} u|^2 + 2n_0 \gamma^2 n(\vect{r}_{\perp},z) |u|^2 \\
        \mathcal{L}_{D}(r;A,W,R,P;\dot{A},\dot{W},\dot{R},\dot{P}) &= n_0 \alpha\,\left[2 |u|^2 \dot{P}-r^2 |u|^2 \dot{R} \right] - r^2 |u|^2 \left(W^4 + R^2\right) + 2 n_0 \gamma^2  n(\vect{r}_{\perp},z) |u|^2\\
    \end{split}
\end{equation}
and then the action becomes
\begin{equation}
    \begin{split}
        J(u,\nabla u) &= \int_{0}^{\Lambda} \iint \mathcal{L}_{D}(u,\nabla u) dx dy \, dz\\   
        J(A,W,R,P;\dot{A},\dot{W},\dot{R},\dot{P}) &= \int_{0}^{\Lambda} n_0 \alpha \left[2 I_0 \dot{P}  - I_2 \dot{R} \right] - I_2 \left(W^4 + R^2\right) +  2n_0 \gamma^2 \left< n(\vect{r}_{\perp},z),|u|^2  \right>   dz\\
    \end{split}
\end{equation}
where 
\begin{equation}
        \left< n(\vect{r}_{\perp},z), |u|^2  \right> =   \frac{A^2W^2}{\pi}  \iint n(\vect{r}_{\perp},z) \exp{-r^2 W^2} dx dy
\end{equation}
is the projection of $|u|^2$ onto the index relatizations $n(\vect{r}_{\perp},z)$ and 
\begin{equation}
    \begin{split}
        I_0 &=  \frac{A^2W^2}{\pi} \int_{0}^{2\pi} \int_{0}^{\infty} \exp{-r^2 W^2} r \,dr \,d\theta =  2A^2W^2 \int_{0}^{\infty} \exp{-r^2 W^2} r \,dr = - A^2 \exp{-\rho^2}\big|_{0}^{\infty} = A^2 
    \end{split}
\end{equation}
\begin{equation}
    \begin{split}
        I_2 &=  \frac{A^2W^2}{\pi} \int_{0}^{2\pi} \int_{0}^{\infty} r^2 \exp{-r^2 W^2} r \,dr \,d\theta = 2 A^2W^2 \int_{0}^{\infty} r^3 \exp{-r^2 W^2} \,dr = \frac{2 A^2W^2}{W^4} \int_{0}^{\infty} \rho^3 \exp{-\rho^2} \,d\rho = \frac{A^2}{W^2} \int_{0}^{\infty} x \exp{-x} \,dx =\frac{A^2}{W^2}       
    \end{split}
\end{equation}
or
\begin{equation}
    \begin{split}
        J(u,\nabla u) &= \int_{0}^{\Lambda} \iint \mathcal{L}_{D}(u,\nabla u) dx dy \, dz\\   
        J(A,W,R,P;\dot{A},\dot{W},\dot{R},\dot{P}) &= \int_{0}^{\Lambda} n_0 \alpha \left[2 I_0 \dot{P}  - I_2 \dot{R} \right] - I_2 \left(W^4 + R^2\right) +  2n_0 \gamma^2 \left< n(\vect{r}_{\perp},z),|u|^2  \right>   dz\\
    \end{split}
\end{equation}
\begin{equation}
    \begin{split}
        J(A,W,R,P;\dot{A},\dot{W},\dot{R},\dot{P}) &= \int_{0}^{\Lambda} n_0 \alpha \left[2A^2 \dot{P} - \frac{A^2 }{W^2}\dot{R} \right]  - \frac{A^2}{W^2} \left(W^4 + R^2\right) +  2n_0 \gamma^2 \left< n(\vect{r}_{\perp},z),|u|^2  \right>  dz\\
    &=\int_{0}^{\Lambda} F_{D}(A,W,R,P;\dot{A},\dot{W},\dot{R},\dot{P})  dz\\
    \end{split}
\end{equation}
where
\begin{equation}
    F_{D}(A,W,R,P;\dot{A},\dot{W},\dot{R},\dot{P}) = n_0 \alpha \left[2A^2 \dot{P} - \frac{A^2 }{W^2}\dot{R} \right]  - \frac{A^2}{W^2} \left(W^4 + R^2\right) +  2n_0 \gamma^2 \left< n(\vect{r}_{\perp},z),|u|^2  \right> 
\end{equation}
 \begin{equation}
    F_{D}(A,W,R,P;\dot{A},\dot{W},\dot{R},\dot{P}) = n_0 \alpha A^2 \left[2 \dot{P} - \frac{1}{W^2}\dot{R} \right]  - A^2 \left(W^2 + \frac{R^2}{W^2}\right) +  2n_0 \gamma^2 \left< n(\vect{r}_{\perp},z),|u|^2  \right> 
\end{equation}

\subsubsection{Euler-Lagrange equation(s)}

Now, we can apply the Euler-Lagrange equations to the reduced density $F_{D}(A,W,R,P;\dot{A},\dot{W},\dot{R},\dot{P})$, i.e.
\begin{equation}
\begin{split}
    \der{F_{D}}{A}  - \der{}{z}\left(\pder{F_{D}}{\dot{A}}\right) = 0  \quad   \der{F_{D}}{W}  - \der{}{z}\left(\pder{F_{D}}{\dot{W}}\right) = 0 \quad  \der{F_{D}}{R}  - \der{}{z}\left(\pder{F_{D}}{\dot{R}}\right) = 0 \quad     \der{F_{D}}{P}  - \der{}{z}\left(\pder{F_{D}}{\dot{P}}\right) = 0. \\
\end{split}
\end{equation}
So lets calculate the terms on the left side of each:
\begin{equation}
\begin{split}
    \der{F_{D}}{A}  =  2n_0 \alpha \left[2\dot{P} - \frac{1}{W^2} \dot{R} \right] A  -  2 \left(W^2 + \frac{R^2}{W^2}\right) A +  2n_0 \gamma^2 \left< n(\vect{r}_{\perp},z),\pder{|u|^2}{A}\right> \quad \quad
    \der{}{z}\left(\pder{F_{D}}{\dot{A}}\right) = \der{}{z}(0) = 0 \\
\end{split}
\end{equation}
\begin{equation}
\begin{split}
    \der{F_{D}}{W}  =  2 n_0 \alpha \frac{A^2 }{W^3} \dot{R}   -  \left(2 A^2 W - 2\frac{A^2 R^2}{W^3}\right)+  2n_0 \gamma^2 \left< n(\vect{r}_{\perp},z),\pder{|u|^2}{W}  \right> \quad \quad 
\der{}{z}\left(\pder{F_{D}}{\dot{W}}\right) = \der{}{z}(0) = 0 
\end{split}
\end{equation}
\begin{equation}
\begin{split}
    \der{F_{D}}{R}  =  -2 \frac{A^2R}{W^2}+  2n_0 \gamma^2 \left< n(\vect{r}_{\perp},z),\pder{|u|^2}{R}  \right> \quad \quad
\der{}{z}\left(\pder{F_{D}}{\dot{R}}\right) &= \der{}{z}\left(  -n_0 \alpha \frac{A^2}{W^2}\right) 
\end{split}
\end{equation}
\begin{equation}
\begin{split}
    \der{F_{D}}{P}  =    2n_0 \gamma^2 \left< n(\vect{r}_{\perp},z),\pder{|u|^2}{P}  \right> \quad \quad
\der{}{z}\left(\pder{F_{D}}{\dot{P}}\right) = \der{}{z}\left( 2 n_0 \alpha A^2 \right)  
\end{split}
\end{equation}


\subsubsection{Equations in P}
Since $|u|^2$ is only dependent on $A,W$, not $R,P$, the last of these equations says
\begin{equation}
\begin{split}
    \der{}{z}\left( 2 n_0 \alpha A^2 \right) = 0,
\end{split}
\end{equation}
or that
\begin{equation}
\begin{split}
    A(z) = A(0)
\end{split}
\end{equation}

\subsubsection{Equations in R}
\begin{equation}
\begin{split}
    \der{F_{D}}{R}  - \der{}{z}\left(\pder{F_{D}}{\dot{R}}\right) &= 0 \\
  -  2 \frac{A^2R}{W^2} + \der{}{z}\left(  n_0 \alpha \frac{A^2}{W^2}\right) &=0 \\
-2 \frac{A^2}{W^2}R + n_0 \alpha \left(  2\frac{A \dot{A}}{W^2} - 2\frac{A^2}{W^3}\dot{W}\right) &=0
\end{split}
\end{equation}
and because $\dot{A} = 0$ we get
\begin{equation}
\begin{split}
 -\frac{A^2}{W^2}R  - n_0 \alpha \frac{A^2}{W^3}\dot{W} &=0 \\
  \dot{W} &= - \frac{1}{n_0 \alpha}\frac{W^3}{ A^2} \frac{A^2R}{W^2} = -\frac{1}{n_0 \alpha} W R   \\
\end{split}
\end{equation}
\begin{equation}
\begin{split}
  \dot{W} = -\frac{1}{n_0 \alpha} W R
\end{split}
\end{equation}

\subsubsection{Equations in W}

\begin{equation}
\begin{split}
    \der{F_{D}}{W}  - \der{}{z}\left(\pder{F_{D}}{\dot{W}}\right) &= 0 \\
    2 n_0 \alpha \frac{A^2 }{W^3} \dot{R}   -  \left(2 A^2 W - 2\frac{A^2 R^2}{W^3}\right)+  2n_0 \gamma^2 \left< n(\vect{r}_{\perp},z),\pder{|u|^2}{W}  \right>   &= 0\\
      2 n_0 \alpha \frac{A^2 }{W^3} \dot{R}   =  2 A^2 W - 2\frac{A^2 R^2}{W^3} -  2n_0 \gamma^2 \left< n(\vect{r}_{\perp},z),\pder{|u|^2}{W}  \right> \\
      \dot{R}  = \frac{1}{n_0 \alpha} W^4 - \frac{1}{n_0 \alpha} R^2  &- \frac{ \gamma^2 }{\alpha} \frac{W^3}{A^2 } \left< n(\vect{r}_{\perp},z),\pder{|u|^2}{W}\right>\\
\end{split}
\end{equation}
\begin{equation}
\begin{split}
      \dot{R}  = \frac{1}{n_0 \alpha} W^4 - \frac{1}{n_0 \alpha} R^2  &- \frac{ \gamma^2 }{\alpha} \frac{W^3}{A^2 } \left< n(\vect{r}_{\perp},z),\pder{|u|^2}{W}\right>\\
\end{split}
\end{equation}

\subsubsection{ Equations in A}
\begin{equation}
\begin{split}
    \der{F_{D}}{A}  - \der{}{z}\left(\pder{F_{D}}{\dot{A}}\right) &= 0 \\
    2n_0 \alpha \left[2\dot{P} - \frac{1}{W^2} \dot{R} \right] A  -  2 \left(W^2 + \frac{R^2}{W^2}\right) A +  2n_0 \gamma^2 \left< n(\vect{r}_{\perp},z),\pder{|u|^2}{A}\right> &= 0 \\
       \left[2\dot{P} - \frac{1}{W^2} \dot{R} \right] A  &=  \frac{2}{2n_0 \alpha} \left(W^2 + \frac{R^2}{W^2}\right) A -  \frac{2n_0 \gamma^2}{2n_0 \alpha} \left< n(\vect{r}_{\perp},z),\pder{|u|^2}{A}\right>\\
      \dot{P}   =  \frac{1}{2 n_0 \alpha}\left(W^2 +  \frac{R^2}{W^2}\right) 
      + \frac{1}{2} \frac{1}{W^2} \dot{R}
      &-  \frac{ \gamma^2 }{2 \alpha}\frac{1}{A}\left< n(\vect{r}_{\perp},z),\pder{|u|^2}{A}\right> \\
\end{split}
\end{equation}
Putting in the equation for $\dot{R}$
\begin{equation}
\begin{split}
      \dot{P}   =  \frac{1}{2 n_0 \alpha}\left(W^2 + \frac{R^2}{W^2}\right) 
      + \frac{1}{2} \frac{1}{W^2}\left[\frac{1}{n_0 \alpha} W^4 - \frac{1}{n_0 \alpha} R^2  - \frac{ \gamma^2 }{\alpha} \frac{W^3}{A^2 } \left< n(\vect{r}_{\perp},z),\pder{|u|^2}{W}\right>\right]
      &-  \frac{ \gamma^2 }{2 \alpha}\frac{1}{A}\left< n(\vect{r}_{\perp},z),\pder{|u|^2}{A}\right> \\
      \dot{P}   =  \frac{1}{n_0 \alpha}W^2
       -  \frac{ \gamma^2 }{2 \alpha} \frac{W}{A^2 } \left<n(\vect{r}_{\perp},z),\pder{|u|^2}{W}\right>
      &-  \frac{ \gamma^2 }{2 \alpha}\frac{1}{A}\left< n(\vect{r}_{\perp},z),\pder{|u|^2}{A}\right> \\      
\end{split}
\end{equation}
\begin{equation}
\begin{split}
      \dot{P}   =  \frac{1}{n_0 \alpha} W^2  
       -  \frac{ \gamma^2 }{2 \alpha} \left<n(\vect{r}_{\perp},z),\frac{W}{A^2} \pder{|u|^2}{W} + \frac{1}{A}\pder{|u|^2}{A} \right>      
\end{split}
\end{equation}

\subsubsection{Finding the modes that project onto the index realizations}

\begin{equation}
\begin{split}
      V_{R} = \frac{W^3}{A^2 } \pder{|u|^2}{W} &= \frac{W^3}{A^2 }\pder{}{W}\left[\frac{A^2W^2}{\pi} \exp{-r^2 W^2}  \right]\\
      &=\frac{W^3}{A^2 }\left[\frac{2 A^2W}{\pi} \exp{-r^2 W^2} - 2 r^2 W\frac{A^2W^2}{\pi} \exp{-r^2 W^2}  \right]\\
      &=\frac{W^3}{A^2 }\left[\frac{2 A^2W}{\pi}  - 2 r^2 \frac{A^2W^3}{\pi} \right]\exp{-r^2 W^2}\\
      &=\frac{W^3}{A^2 }\frac{2 A^2W}{\pi} \left[1  -  r^2 W^2\right]\exp{-r^2 W^2}\\
      &=\frac{2 W^4}{\pi}\left[1  -  r^2 W^2\right]\exp{-r^2 W^2}\\
\end{split}
\end{equation}

\begin{equation}
\begin{split}
      V_{P} = \frac{W}{A^2 } \pder{|u|^2}{W} + \frac{1}{A}\pder{|u|^2}{A} &= \frac{W}{A^2}\pder{}{W}\left[\frac{A^2W^2}{\pi} \exp{-r^2 W^2}\right] + \frac{1}{A}\pder{}{A}\left[\frac{A^2W^2}{\pi} \exp{-r^2 W^2}\right] \\
      \frac{W}{A^2 } \pder{|u|^2}{W} +\frac{1}{A}\pder{|u|^2}{A} &= \frac{W}{A^2}\left[\frac{2 A^2W}{\pi} \exp{-r^2 W^2} - 2 r^2 W\frac{A^2W^2}{\pi} \exp{-r^2 W^2}  \right] + \frac{1}{A}\left[2\frac{AW^2}{\pi} \exp{-r^2 W^2}\right] \\
    \frac{W}{A^2 } \pder{|u|^2}{W} +\frac{1}{A}\pder{|u|^2}{A} &=\left[ \frac{W}{A^2}\frac{2 A^2W}{\pi} \exp{-r^2 W^2} - 2 r^2  \frac{W}{A^2}W\frac{A^2W^2}{\pi} \exp{-r^2 W^2}  \right] + \left[2\frac{W^2}{\pi} \exp{-r^2 W^2}\right] \\
    \frac{W}{A^2 } \pder{|u|^2}{W} +\frac{1}{A}\pder{|u|^2}{A} &=\left[\frac{2 W^2}{\pi} \exp{-r^2 W^2} - 2 r^2 \frac{W^4}{\pi} \exp{-r^2 W^2}  \right] + \left[2\frac{W^2}{\pi} \exp{-r^2 W^2}\right] \\   
    \frac{W}{A^2 } \pder{|u|^2}{W} +\frac{1}{A}\pder{|u|^2}{A} &=2\frac{W^2}{\pi} \left[ 2 - r^2 W^2\right] \exp{-r^2 W^2}
\end{split}
\end{equation}

\subsubsection{Calculating statistics}
The two stochastic forcing terms are writen as the projections onto the index process. Since projections are linear and the index process is typically modeled as a mean zero Gaussian process, the resulting projections will also be mean Gaussian processes in the $z$ variable.  So we only need to find the autocovariance and covariance of the processes to simulate them.  First, lets name them:
\begin{equation}
\begin{split}
      N_{R}(z)  = \left< n(\vect{r}_{\perp},z),V_{R}(r,z)\right> = \iint n(\vect{r}_{\perp},z)V_{R}(r,z)dx dy\\
\end{split}
\end{equation}
where 
\begin{equation}
\begin{split}
      V_{R} = \frac{2}{\pi} W^4 \left[1  -  r^2 W^2\right]\exp{-r^2 W^2}\\
\end{split}
\end{equation}
and
\begin{equation}
\begin{split}
      N_{P}(z)  = \left< n(\vect{r}_{\perp},z),V_{P}(r,z)\right> =  \iint n(\vect{r}_{\perp},z)V_{P}(r,z) dx dy\\
\end{split}
\end{equation}
where 
\begin{equation}
\begin{split}
      V_{P} = \frac{2}{\pi} W^2 \left[ 2 - r^2 W^2\right] \exp{-r^2 W^2}\\ \\
\end{split}
\end{equation}


\subsubsection{Calculating the auto-covariances}
\begin{equation}
\begin{split}
      \Expect{N_{R}(z_1)N_{R}(z_2)}  =  \iint \iint \Expect{n(x_1,y_1,z_1)n(x_2,y_2,z_2)} V_{R}(x_1,y_1,z_1)V_{R}(x_2,y_2,z_2) dx_1 dy_1 dx_2 dy_2\\
\end{split}
\end{equation}




\section{Gaussian Ansatz: General}

We begin by using a Gaussian ansatz:
\begin{equation}
    u(\vect{r}_{\perp},\vect{p}(z)) = I(\vect{p}) \exp{- \left(\Theta(\vect{r}_{\perp},\vect{p}) + i \Phi(\vect{r}_{\perp},\vect{p}) \right)}
\end{equation}
Now, lets calculate each term in the Lagrangian density:
\begin{equation}
    \begin{split}
    \partial_{z}(u) &= \left[ \frac{\nabla_{\vect{p}}{I}\cdot\vect{\dot{p}}}{I} - \nabla_{\vect{p}}{\Theta}\cdot\vect{\dot{p}} - i \nabla_{\vect{p}}{\Phi}\cdot\vect{\dot{p}} \right] u\\
    \partial_{z}(u)\conj{u} &= \left[ \frac{\nabla_{\vect{p}}{I}\cdot\vect{\dot{p}}}{I} -\nabla_{\vect{p}}{\Theta}\cdot\vect{\dot{p}} - i \nabla_{\vect{p}}{\Phi}\cdot\vect{\dot{p}} \right] |u|^2 \\
   -2 n_0 \alpha \Im{\partial_{z}(u)\conj{u}} &= 2 n_0 \alpha \left(\nabla_{\vect{p}}{\Phi}\cdot\vect{\dot{p}}\right) |u|^2 \\    
   &= 2 n_0 \alpha \left(\nabla_{\vect{p}}{\Phi}\cdot\vect{\dot{p}}\right) I^2(\vect{p})\exp{-2 \Theta(\vect{r}_{\perp},\vect{p})} \\    
    \end{split}
\end{equation}
\begin{equation}
    \begin{split}
        \nabla_{\perp} u &= -\left(\nabla_{\perp}\Theta(\vect{r}_{\perp},\vect{p}) + i \nabla_{\perp}\Phi(\vect{r}_{\perp},\vect{p}) \right) u \\
    -|\nabla_{\perp} u|^2 &= - \left(|\nabla_{\perp}\Theta(\vect{r}_{\perp},\vect{p})|^2 + |\nabla_{\perp}\Phi(\vect{r}_{\perp},\vect{p})|^2 \right)|u|^2 \\
    &= - \left(|\nabla_{\perp}\Theta(\vect{r}_{\perp},\vect{p})|^2 + |\nabla_{\perp}\Phi(\vect{r}_{\perp},\vect{p})|^2 \right)I^2(\vect{p})\exp{-2 \Theta(\vect{r}_{\perp},\vect{p})} 
    \end{split}
\end{equation}
\begin{equation}
    \begin{split}
    |u|^2 &= I^2(\vect{p})\exp{-2 \Theta(\vect{r}_{\perp},\vect{p})}  
    \end{split}
\end{equation}

\begin{equation}
\begin{split}
    \mathcal{L}_{D}(u,\nabla u) &= -2 n_0 \alpha\, \mathrm{Im}\left(\partial_{z}(u) \overline{u}\right) - |\nabla_{\perp} u|^2 + 2n_0 \gamma^2 n(\vect{r}_{\perp},z) |u|^2 \\
    & = \left(2 n_0 \alpha \nabla_{\vect{p}}{\Phi}\cdot\vect{\dot{p}} - \left(|\nabla_{\perp}\Theta(\vect{r}_{\perp},\vect{p})|^2 + |\nabla_{\perp}\Phi(\vect{r}_{\perp},\vect{p})|^2 \right) + 2n_0 \gamma^2 n(\vect{r}_{\perp},z)\right) I^2(\vect{p})\exp{-2 \Theta(\vect{r}_{\perp},\vect{p})}
\end{split}
\end{equation}

\subsection{Tilted Gaussian}
\begin{equation}
    \vect{p} = \left[A, W_x,W_y; T_x, T_y, X, Y; F_x,F_y, P\right]
\end{equation}
\begin{equation}
    \begin{split}
        I(\vect{p}) &= \frac{A \sqrt{W_x W_y}}{\sqrt{\pi}} \\
        \Theta(\vect{r}_{\perp},\vect{p}) &= \frac{1}{2} \left( W_x^2(x-X)^2 + W_y^2(y-Y)^2 \right) \\
        \Phi(\vect{r}_{\perp},\vect{p}) &= P + T_x (x-X) + T_y (y-Y) + F_x (x-X)^2 + F_{y} (y-Y)^2
    \end{split}
\end{equation}




\begin{equation}
    \vect{p} = \left[A, W_x,W_y; T_x, T_y, X, Y; F_x,F_y, P\right]
\end{equation}
\begin{equation}
    \begin{split}
        \nabla_{\vect{p}}{\Phi}\cdot\vect{\dot{p}} &= \dot{P} + \dot{T}_x (x-X) + \dot{T}_y (y-Y) + \dot{F}_{x} (x-X)^2 + \dot{F}_{y} (y-Y)^2 \\
        &\quad - \dot{X} (T_x + 2F_x(x-X)) - \dot{Y} (T_y + 2F_y(y-Y)) \\
    \end{split}
\end{equation}

\begin{equation}
    \begin{split}
        \nabla_{\perp}\Theta(\vect{r}_{\perp},\vect{p}) &= \left[ W_x^2 (x-X),  W_y^2 (y-Y) \right] \\
        |\nabla_{\perp}\Theta(\vect{r}_{\perp},\vect{p})|^2 &= W_x^4 (x-X)^2 +  W_y^4 (y-Y)^2 \\
    \end{split}
\end{equation}

\begin{equation}
    \begin{split}
        \nabla_{\perp}\Phi(\vect{r}_{\perp},\vect{p}) &= \left[T_x + 2 F_x (x-X), T_y + 2 F_{y} (y-Y)\right] \\
        |\nabla_{\perp}\Phi(\vect{r}_{\perp},\vect{p})|^2 &= \left(T_x + 2 F_x (x-X)\right)^2 +  \left(T_y + 2 F_{y} (y-Y)\right)^2
    \end{split}
\end{equation}

\begin{equation}
    \begin{split}
    \mathcal{L}_{D}(u,\nabla u) & = 2 n_0 \alpha\, \left(\dot{P} + \dot{T}_x (x-X) + \dot{T}_y (y-Y) + \dot{F}_{x} (x-X)^2 + \dot{F}_{y} (y-Y)^2\right)\frac{A^2W_x W_y}{\pi} \exp{-\left( W_x^2(x-X)^2 + W_y^2(y-Y)^2 \right)} \\
        &\quad - \dot{X} (T_x + 2F_x(x-X)) \frac{A^2W_x W_y}{\pi} \exp{-\left( W_x^2(x-X)^2 + W_y^2(y-Y)^2 \right)}
\\
&- \dot{Y} (T_y + 2F_y(y-Y)) \frac{A^2W_x W_y}{\pi} \exp{-\left( W_x^2(x-X)^2 + W_y^2(y-Y)^2 \right)}
\\ 
        &- \left(W_x^4 (x-X)^2 +  W_y^4 (y-Y)^2\right) \frac{A^2W_x W_y}{\pi} \exp{-\left( W_x^2(x-X)^2 + W_y^2(y-Y)^2 \right)}
\\
        &- \left(T_x + 2 F_x (x-X)\right)^2 \frac{A^2W_x W_y}{\pi} \exp{-\left( W_x^2(x-X)^2 + W_y^2(y-Y)^2 \right)}
\\
        &-  \left(T_y + 2 F_{y} (y-Y)\right)^2 \frac{A^2W_x W_y}{\pi} \exp{-\left( W_x^2(x-X)^2 + W_y^2(y-Y)^2 \right)}
\\
        &+ 2n_0 \gamma^2 n(\vect{r}_{\perp},z) \frac{A^2W_x W_y}{\pi} \exp{-\left( W_x^2(x-X)^2 + W_y^2(y-Y)^2 \right)} \\
    \end{split}
\end{equation}
\begin{equation}
    \begin{split}
 \mathcal{F}_{D}(\vect{p},\dot{\vect{p}}) &= 2 n_0 \alpha\, \frac{A^2W_x W_y}{\pi} \left(\dot{P} I_0 + \dot{T}_x I_{1,x}  + \dot{T}_y I_{1,y} + \dot{F}_{x} I_{2,x} + \dot{F}_{y} I_{2,x}\right) \\
        &\quad - \dot{X} \frac{A^2W_x W_y}{\pi} (T_x I_0 + 2 F_x I_{1,x}) \\
&- \dot{Y}\frac{A^2W_x W_y}{\pi} (T_y I_0 + 2 F_y I_{1,y}) \\ 
        &- \frac{A^2W_x W_y}{\pi}\left(W_x^4 I_{2,x} +  W_y^4 I_{2,y}\right) \\
        &- \frac{A^2W_x W_y}{\pi}\left(T_x^2 I_0 + 2 T_{x} F_{x} I_{1,x} + 4 F_{x}^2 I_{2,x}\right) \\
        &- \frac{A^2W_x W_y}{\pi}\left(T_y^2 I_0 + 2 T_{y} F_{y} I_{1,y} + 4 F_{y}^2 I_{2,y}\right) \\
        &+ 2n_0 \gamma^2 \left<n(\vect{r}_{\perp},z), \frac{A^2W_x W_y}{\pi} \exp{-\left( W_x^2(x-X)^2 + W_y^2(y-Y)^2 \right)}\right>
    \end{split}
\end{equation}


\begin{equation}
    \begin{split}
       I_0 &= \iint^{\infty}_{\infty} \exp{-\left( W_x^2(x-X)^2 + W_y^2(y-Y)^2 \right)} \, dx \,dy =   \frac{\pi}{W_{x}W_{y}}\\
       I_{1,x} &= \iint^{\infty}_{\infty} (x-X) \exp{-\left( W_x^2(x-X)^2 + W_y^2(y-Y)^2 \right)} \, dx \,dy = 0\\
       I_{1,y} &= \iint^{\infty}_{\infty} (y-Y) \exp{-\left( W_x^2(x-X)^2 + W_y^2(y-Y)^2 \right)} \, dx \,dy = 0\\
       I_{2,x} &= \iint^{\infty}_{\infty} (x-X)^2 \exp{-\left( W_x^2(x-X)^2 + W_y^2(y-Y)^2 \right)} \, dx \,dy = \frac{\pi}{2W_{x}^3W_{y}}\\
       I_{2,y} &= \iint^{\infty}_{\infty} (y-Y)^2 \exp{-\left( W_x^2(x-X)^2 + W_y^2(y-Y)^2 \right)} \, dx \,dy = \frac{\pi}{2W_{y}^3W_{x}}\\
    \end{split}
\end{equation}


\begin{equation}
    \begin{split}
    \mathcal{F}_{D}(\vect{p},\dot{\vect{p}}) & = 2 n_0 \alpha\, \frac{A^2W_x W_y}{\pi} \left(\dot{P} \frac{\pi}{W_{x}W_{y}}  + \dot{F}_{x} \frac{\pi}{2W_{x}^3W_{y}} + \dot{F}_{y} \frac{\pi}{2W_{y}^3W_{x}} \right) \\
        &\quad - \dot{X} \frac{A^2W_x W_y}{\pi} T_x \frac{\pi}{W_{x}W_{y}} \\
&- \dot{Y}\frac{A^2W_x W_y}{\pi} T_y \frac{\pi}{W_{x}W_{y}} \\ 
        &- \frac{A^2W_x W_y}{\pi}\left(W_x^4 \frac{\pi}{2W_{x}^3W_{y}} +  W_y^4 \frac{\pi}{2W_{y}^3W_{x}}\right) \\
        &- \frac{A^2W_x W_y}{\pi}\left(T_x^2 \frac{\pi}{W_{x}W_{y}}  + 4 F_{x}^2 \frac{\pi}{2W_{x}^3W_{y}}\right) \\
        &- \frac{A^2W_x W_y}{\pi}\left(T_y^2 \frac{\pi}{W_{x}W_{y}}  + 4 F_{y}^2 \frac{\pi}{2W_{y}^3W_{x}}\right) \\
        &+ 2n_0 \gamma^2 \left<n(\vect{r}_{\perp},z), \frac{A^2W_x W_y}{\pi} \exp{-\left( W_x^2(x-X)^2 + W_y^2(y-Y)^2 \right)}\right>
    \end{split}
\end{equation}

\begin{equation}
    \begin{split}
    \mathcal{F}_{D}(\vect{p},\dot{\vect{p}}) & = 2 n_0 \alpha\, A^2 \left(\dot{P}  + \frac{\dot{F}_{x} }{2W_{x}^2} + \frac{\dot{F}_{y}}{2W_{y}^2} \right) - A^2 \dot{X}  T_x - A^2  \dot{Y} T_y  - A^2 \left( \frac{W_x^2}{2} +   \frac{W_y^2}{2}\right) \\
        &- A^2 \left(T_x^2 + 2  \frac{F_{x}^2}{W_{x}^2}\right) - A^2 \left(T_y^2 + 2 \frac{F_{y}^2}{W_{y}^2}\right) \\
        &+ 2n_0 \gamma^2 \left<n(\vect{r}_{\perp},z), \frac{A^2W_x W_y}{\pi} \exp{-\left( W_x^2(x-X)^2 + W_y^2(y-Y)^2 \right)}\right>
    \end{split}
\end{equation}


\begin{equation}
    \begin{split}
    \mathcal{F}_{D}(\vect{p},\dot{\vect{p}}) & = 2 n_0 \alpha\, A^2 \left(\dot{P}  + \frac{\dot{F}_{x} }{2W_{x}^2} + \frac{\dot{F}_{y}}{2W_{y}^2} \right) - A^2 \left(\dot{X}T_x + \dot{Y} T_y\right)  - \frac{A^2}{2} \left(W_x^2+ W_y^2\right) \\
        &- A^2 \left(T_x^2 + T_y^2 \right) - 2 A^2 \left(\frac{F_{x}^2}{W_{x}^2} + \frac{F_{y}^2}{W_{y}^2}\right) \\
        &+ 2n_0 \gamma^2 \left<n(\vect{r}_{\perp},z), \frac{A^2W_x W_y}{\pi} \exp{-\left( W_x^2(x-X)^2 + W_y^2(y-Y)^2 \right)}\right>
    \end{split}
\end{equation}


\subsection{Extracting Parameters}

\begin{equation}
\begin{split}
\int |u|^2 \,dx\,dy = \frac{A^2W_x W_y}{\pi} \int \exp{-\left( W_x^2(x-X)^2 + W_y^2(y-Y)^2 \right)}\,dx\,dy = \\
\frac{A^2}{\pi}\int \exp{-\left( \hat{x}^2 + \hat{y}^2 \right)}\,d\hat{x}\,d\hat{y} = A^2 
\end{split}
\end{equation}


\begin{equation}
\begin{split}
\int x |u|^2 \,dx\,dy = \frac{A^2W_x W_y}{\pi} \int x \exp{-\left( W_x^2(x-X)^2 + W_y^2(y-Y)^2 \right)}\,dx\,dy = \\
\frac{A^2}{\sqrt{\pi}} \int (\frac{\hat{x}}{W_x} + X) \exp{-\hat{x}^2 }\,d\hat{x} = A^2 = A^2 X
\end{split}
\end{equation}

\begin{equation}
\begin{split}
\int (x-X)^2 |u|^2 \,dx\,dy = \frac{A^2W_x W_y}{\pi} \int  (x-X)^2 \exp{-\left( W_x^2(x-X)^2 + W_y^2(y-Y)^2 \right)}\,dx\,dy = \\
\frac{A^2}{\sqrt{\pi}} \int \frac{\hat{x}^2}{W_x^2} \exp{-\hat{x}^2 }\,d\hat{x} = \frac{A^2}{2 W_x^2}
\end{split}
\end{equation}


\subsection{Euler-Lagrange equation(s)}
Now, we can apply the Euler-Lagrange equations to the reduced density $\mathcal{F}_{D}(\vect{p},\dot{\vect{p}})$, i.e.
\begin{equation}
\begin{split}
    \der{\mathcal{F}_{D}}{p_i}  - \der{}{z}\left(\pder{\mathcal{F}_{D}}{\dot{p_i}}\right) = 0 
\end{split}
\end{equation}
for $p_i = A, W_x,W_y; T_x, T_y, X, Y; F_x,F_y, P$.
So lets calculate the terms on the left side of each:
\begin{equation}
\begin{split}
    \der{F_{D}}{A}  &=  2 n_0 \alpha\, (2A) \left(\dot{P}  + \frac{\dot{F}_{x} }{2W_{x}^2} + \frac{\dot{F}_{y}}{2W_{y}^2} \right) - 2A  \left(\dot{X}T_x + \dot{Y} T_y\right)  - A  \left(W_x^2+ W_y^2\right) \\
        &- 2A  \left(T_x^2 + T_y^2 \right) - 4A  \left(\frac{F_{x}^2}{W_{x}^2} + \frac{F_{y}^2}{W_{y}^2}\right) +  2n_0 \gamma^2 \left< n(\vect{r}_{\perp},z),\pder{|u|^2}{A}\right>\\
    \der{}{z}\left(\pder{F_{D}}{\dot{A}}\right) &= \der{}{z}(0) = 0 \\
\end{split}
\end{equation}

\begin{equation}
\begin{split}
    \der{F_{D}}{W_x} & = 2 n_0 \alpha\, A^2 \left(  -\frac{\dot{F}_{x} }{W_{x}^3} \right)  - A^2 W_x + 4 A^2 \frac{F_{x}^2}{W_{x}^3}  + 2n_0 \gamma^2 \left<n(\vect{r}_{\perp},z), \pder{|u|^2}{W_{x}}\right>\\
\der{}{z}\left(\pder{F_{D}}{\dot{W_x}}\right) &= \der{}{z}(0) = 0 
\end{split}
\end{equation}

\begin{equation}
\begin{split}
    \der{F_{D}}{W_y} & = 2 n_0 \alpha\, A^2 \left(  -\frac{\dot{F}_{y} }{W_{y}^3} \right)  - A^2 W_y + 4 A^2 \frac{F_{y}^2}{W_{y}^3}  + 2n_0 \gamma^2 \left<n(\vect{r}_{\perp},z), \pder{|u|^2}{W_{y}}\right>\\
\der{}{z}\left(\pder{F_{D}}{\dot{W_y}}\right) &= \der{}{z}(0) = 0 
\end{split}
\end{equation}

\begin{equation}
\begin{split}
    \der{F_{D}}{T_x} & = -A^2 \dot{X} - 2A^2 T_x  + 2n_0 \gamma^2 \left<n(\vect{r}_{\perp},z), \pder{|u|^2}{T_{x}}\right>\\
\der{}{z}\left(\pder{F_{D}}{\dot{T_x}}\right) &= \der{}{z}(0) = 0 
\end{split}
\end{equation}

\begin{equation}
\begin{split}
    \der{F_{D}}{T_y} & = -A^2 \dot{Y} - 2A^2 T_y  + 2n_0 \gamma^2 \left<n(\vect{r}_{\perp},z), \pder{|u|^2}{T_{y}}\right>\\
\der{}{z}\left(\pder{F_{D}}{\dot{T_y}}\right) &= \der{}{z}(0) = 0 
\end{split}
\end{equation}


\begin{equation}
\begin{split}
    \der{F_{D}}{X} & = 2n_0 \gamma^2 \left<n(\vect{r}_{\perp},z), \pder{|u|^2}{X}\right>\\
\der{}{z}\left(\pder{F_{D}}{\dot{X}}\right) &= \der{}{z}\left(-A^2 T_x \right)
\end{split}
\end{equation}

\begin{equation}
\begin{split}
    \der{F_{D}}{Y} & = 2n_0 \gamma^2 \left<n(\vect{r}_{\perp},z), \pder{|u|^2}{Y}\right>\\
\der{}{z}\left(\pder{F_{D}}{\dot{Y}}\right) &= \der{}{z}\left(-A^2 T_y \right)
\end{split}
\end{equation}


\begin{equation}
\begin{split}
    \der{F_{D}}{F_x}  &=  - 4 A^2 \frac{F_{x}}{W_{x}^2} + 2n_0 \gamma^2 \left<n(\vect{r}_{\perp},z),\pder{|u|^2}{F_x}\right> \\
    \der{}{z}\left(\pder{F_{D}}{\dot{F_x}}\right) &= 2 n_0 \alpha \der{}{z}\left(A^2 \frac{1}{2W_{x}^2} \right)
\end{split}
\end{equation}

\begin{equation}
\begin{split}
    \der{F_{D}}{F_y}  &=  - 4 A^2 \frac{F_{y}}{W_{y}^2} + 2n_0 \gamma^2 \left<n(\vect{r}_{\perp},z),\pder{|u|^2}{F_y}\right> \\
    \der{}{z}\left(\pder{F_{D}}{\dot{F_y}}\right) &= 2 n_0 \alpha  \der{}{z}\left(A^2 \frac{1}{2W_{y}^2} \right)
\end{split}
\end{equation}

\begin{equation}
\begin{split}
    \der{F_{D}}{P}  &= 2n_0 \gamma^2 \left< n(\vect{r}_{\perp},z),\pder{|u|^2}{P}  \right> \\
\der{}{z}\left(\pder{F_{D}}{\dot{P}}\right) &= 2 n_0 \alpha \der{}{z}\left(A^2 \right)  
\end{split}
\end{equation}


\subsubsection{Equations in P}
Since $|u|^2$ is only dependent on $A,W$, not $R,P$, the last of these equations says
\begin{equation}
\begin{split}
    \der{}{z}\left( 2 n_0 \alpha A^2 \right) = 0,
\end{split}
\end{equation}
or that
\begin{equation}
\begin{split}
    A(z) = A(0)
\end{split}
\end{equation}


\subsubsection{Equations in $F_x$, $F_y$}

\begin{equation}
\begin{split}
    - 4 A^2 \frac{F_{x}}{W_{x}^2} + 2n_0 \gamma^2 \left<n(\vect{r}_{\perp},z),\pder{|u|^2}{F_x}\right> &= 2 n_0 \alpha \der{}{z}\left(A^2 \frac{1}{2W_{x}^2} \right) = -2 n_0 \alpha A^2 \frac{\dot{W_{x}}}{W_{x}^3} \\
      A^2 \frac{\dot{W_{x}}}{W_{x}^3} &= \frac{4}{2 n_0 \alpha} A^2 \frac{F_{x}}{W_{x}^2} - \frac{2n_0 \gamma^2}{2 n_0 \alpha}  \left<n(\vect{r}_{\perp},z),\pder{|u|^2}{F_x}\right>\\
     \dot{W_{x}} &= \frac{2}{n_0 \alpha} F_{x}W_{x} - \frac{\gamma^2}{\alpha} \frac{W_{x}^3}{A^2} \left<n(\vect{r}_{\perp},z),\pder{|u|^2}{F_x}\right> \\
     \dot{W_{x}} &= \frac{2}{n_0 \alpha} F_{x}W_{x}
    \end{split}
\end{equation}

\begin{equation}
\begin{split}
  - 4 A^2 \frac{F_{y}}{W_{y}^2} + 2n_0 \gamma^2 \left<n(\vect{r}_{\perp},z),\pder{|u|^2}{F_y}\right>
    &= 2 n_0 \alpha \der{}{z}\left(A^2 \frac{1}{2W_{y}^2} \right) = -2 n_0 \alpha A^2 \frac{\dot{W_{y}}}{2W_{y}^3} \\
   \dot{W_{y}} &= \frac{2}{n_0 \alpha} F_{y}W_{y} - \frac{\gamma^2}{\alpha} \frac{W_{y}^3}{A^2} \left<n(\vect{r}_{\perp},z),\pder{|u|^2}{F_y}\right> \\
      \dot{W_{y}} &= \frac{2}{n_0 \alpha} F_{y}W_{y} 
\end{split}
\end{equation}


\subsubsection{Equations in $X$,$Y$}

\begin{equation}
\begin{split}
    2n_0 \gamma^2 \left<n(\vect{r}_{\perp},z), \pder{|u|^2}{X}\right> &= \der{}{z}\left(-A^2 T_x \right) =  -A^2 \dot{T_x}\\
        \dot{T_x} &= - 2n_0 \gamma^2 \frac{1}{A^2} \left<n(\vect{r}_{\perp},z), \pder{|u|^2}{X}\right> 
\end{split}
\end{equation}

\begin{equation}
\begin{split}
   2n_0 \gamma^2 \left<n(\vect{r}_{\perp},z), \pder{|u|^2}{Y}\right> &= \der{}{z}\left(-A^2 T_y \right) \\
        \dot{T_y} &= - 2n_0 \gamma^2 \frac{1}{A^2} \left<n(\vect{r}_{\perp},z), \pder{|u|^2}{Y}\right> 
\end{split}
\end{equation}

\subsubsection{Equations in $T_x$,$T_y$}

\begin{equation}
\begin{split}
   -A^2 \dot{X} - 2A^2 T_x  + 2n_0 \gamma^2 \left<n(\vect{r}_{\perp},z), \pder{|u|^2}{T_{x}}\right> &= 0\\
   \dot{X} &= - 2T_x  + 2n_0 \gamma^2 \frac{1}{A^2} \left<n(\vect{r}_{\perp},z), \pder{|u|^2}{T_{x}}\right>\\
   \dot{X} &= - 2T_x 
\end{split}
\end{equation}

\begin{equation}
\begin{split}
   -A^2 \dot{Y} - 2A^2 T_y  + 2n_0 \gamma^2 \left<n(\vect{r}_{\perp},z), \pder{|u|^2}{T_{y}}\right> &=0\\
      \dot{Y} &= - 2T_y  + 2n_0 \gamma^2 \frac{1}{A^2} \left<n(\vect{r}_{\perp},z), \pder{|u|^2}{T_{y}}\right>\\
      \dot{Y} &= - 2T_y 
\end{split}
\end{equation}

\subsubsection{Equations in $W_x$,$W_y$}

\begin{equation}
\begin{split}
     2 n_0 \alpha\, A^2 \left(  -\frac{\dot{F}_{x} }{W_{x}^3} \right)  &- A^2 W_x + 4 A^2 \frac{F_{x}^2}{W_{x}^3}  + 2n_0 \gamma^2 \left<n(\vect{r}_{\perp},z), \pder{|u|^2}{W_{x}}\right> = 0 \\
     2 n_0 \alpha\, A^2 \frac{\dot{F}_{x} }{W_{x}^3}  &=   - A^2 W_x + 4 A^2 \frac{F_{x}^2}{W_{x}^3}  + 2n_0 \gamma^2 \left<n(\vect{r}_{\perp},z), \pder{|u|^2}{W_{x}}\right>  \\
     \frac{\dot{F}_{x} }{W_{x}^3}  &=   - \frac{A^2 W_x}{2 n_0 \alpha\, A^2} + \frac{4 A^2 }{2 n_0 \alpha\, A^2} \frac{F_{x}^2}{W_{x}^3}  + \frac{2n_0 \gamma^2}{2 n_0 \alpha\, A^2}  \left<n(\vect{r}_{\perp},z), \pder{|u|^2}{W_{x}}\right> \\
     \dot{F}_{x}  &=   - \frac{W_{x}^4}{2 n_0 \alpha} + \frac{2 W_{x}^3}{n_0 \alpha} \frac{F_{x}^2}{W_{x}^3}  + \frac{\gamma^2 W_{x}^3}{ \alpha\, A^2}  \left<n(\vect{r}_{\perp},z), \pder{|u|^2}{W_{x}}\right> \\ 
     \dot{F}_{x}  &=   - \frac{1}{2 n_0 \alpha} W_{x}^4+ \frac{2}{n_0 \alpha} F_{x}^2  + \frac{\gamma^2 W_{x}^3}{ \alpha\, A^2}  \left<n(\vect{r}_{\perp},z), \pder{|u|^2}{W_{x}}\right> \\    
\end{split}
\end{equation}

\begin{equation}
\begin{split}
    2 n_0 \alpha\, A^2 \left(  -\frac{\dot{F}_{y} }{W_{y}^3} \right) & - A^2 W_y + 4 A^2 \frac{F_{y}^2}{W_{y}^3}  + 2n_0 \gamma^2 \left<n(\vect{r}_{\perp},z), \pder{|u|^2}{W_{y}}\right>= 0 \\
         \dot{F}_{y}  &=   - \frac{1}{2 n_0 \alpha} W_{y}^4 + \frac{2}{n_0 \alpha} F_{y}^2  + \frac{\gamma^2 W_{y}^3}{ \alpha\, A^2}  \left<n(\vect{r}_{\perp},z), \pder{|u|^2}{W_{y}}\right> \\
\end{split}
\end{equation}

\subsubsection{ Equations in A}
\begin{equation}
\begin{split}
    &2 n_0 \alpha\, (2A) \left(\dot{P}  + \frac{\dot{F}_{x} }{2W_{x}^2} + \frac{\dot{F}_{y}}{2W_{y}^2} \right) - 2A  \left(\dot{X}T_x + \dot{Y} T_y\right)  - A  \left(W_x^2+ W_y^2\right) \\
    &- 2A  \left(T_x^2 + T_y^2 \right) - 4A  \left(\frac{F_{x}^2}{W_{x}^2} + \frac{F_{y}^2}{W_{y}^2}\right) +  2n_0 \gamma^2 \left< n(\vect{r}_{\perp},z),\pder{|u|^2}{A}\right> = 0\\
    & \left(\dot{P}  + \frac{\dot{F}_{x} }{2W_{x}^2} + \frac{\dot{F}_{y}}{2W_{y}^2} \right) = \frac{2A}{2 n_0 \alpha\, (2A)}  \left(\dot{X}T_x + \dot{Y} T_y\right)  + \frac{A}{2 n_0 \alpha\, (2A)}  \left(W_x^2+ W_y^2\right) \\
    &\quad+ \frac{2A}{2 n_0 \alpha\, (2A)}  \left(T_x^2 + T_y^2 \right) + \frac{4A}{2 n_0 \alpha\, (2A)}  \left(\frac{F_{x}^2}{W_{x}^2} + \frac{F_{y}^2}{W_{y}^2}\right) -  \frac{2n_0 \gamma^2}{2 n_0 \alpha\, (2A)} \left< n(\vect{r}_{\perp},z),\pder{|u|^2}{A}\right> \\ 
    & \left(\dot{P}  + \frac{\dot{F}_{x} }{2W_{x}^2} + \frac{\dot{F}_{y}}{2W_{y}^2} \right) = \frac{1}{2 n_0 \alpha}  \left(\dot{X}T_x + \dot{Y} T_y\right)  + \frac{1}{4 n_0 \alpha}  \left(W_x^2+ W_y^2\right) \\
    &\quad + \frac{1}{2 n_0 \alpha}  \left(T_x^2 + T_y^2 \right) + \frac{1}{n_0 \alpha}  \left(\frac{F_{x}^2}{W_{x}^2} + \frac{F_{y}^2}{W_{y}^2}\right) -  \frac{\gamma^2}{2 \alpha A)} \left< n(\vect{r}_{\perp},z),\pder{|u|^2}{A}\right>    
\end{split}
\end{equation}

\begin{equation}
\begin{split}
    \dot{P}  &= -\frac{\dot{F}_{x} }{2W_{x}^2} - \frac{\dot{F}_{y}}{2W_{y}^2} 
    + \frac{1}{2 n_0 \alpha}  \left(\dot{X}T_x + \dot{Y} T_y\right) 
     + \frac{1}{4 n_0 \alpha}  \left(W_x^2+ W_y^2\right) \\    
    &\quad + \frac{1}{2 n_0 \alpha}  \left(T_x^2 + T_y^2 \right) + \frac{1}{n_0 \alpha}  \left(\frac{F_{x}^2}{W_{x}^2} + \frac{F_{y}^2}{W_{y}^2}\right) -  \frac{\gamma^2}{2 \alpha A} \left< n(\vect{r}_{\perp},z),\pder{|u|^2}{A}\right>    
\end{split}
\end{equation}



\begin{equation}
\begin{split}
    \dot{P}  &= -\frac{1}{2W_{x}^2} \left(- \frac{1}{2 n_0 \alpha} W_{x}^4+ \frac{2}{n_0 \alpha} F_{x}^2  + \frac{\gamma^2 W_{x}^3}{ \alpha\, A^2}  \left<n(\vect{r}_{\perp},z), \pder{|u|^2}{W_{x}}\right>\right)  \\
    &\quad - \frac{1}{2W_{y}^2} \left(- \frac{1}{2 n_0 \alpha} W_{y}^4 + \frac{2}{n_0 \alpha} F_{y}^2  + \frac{\gamma^2 W_{y}^3}{ \alpha\, A^2}  \left<n(\vect{r}_{\perp},z), \pder{|u|^2}{W_{y}}\right>\right)\\
    &\quad + \frac{1}{2 n_0 \alpha}T_x  \left( - 2T_x \right) \\
    &\quad + \frac{1}{2 n_0 \alpha}T_y  \left( - 2T_y \right) \\
    &\quad + \frac{1}{4 n_0 \alpha}  \left(W_x^2+ W_y^2\right) \\    
    &\quad + \frac{1}{2 n_0 \alpha}  \left(T_x^2 + T_y^2 \right) \\
    &\quad + \frac{1}{n_0 \alpha}  \left(\frac{F_{x}^2}{W_{x}^2} + \frac{F_{y}^2}{W_{y}^2}\right) \\
    &\quad -  \frac{\gamma^2}{2 \alpha A} \left< n(\vect{r}_{\perp},z),\pder{|u|^2}{A}\right>   \\
    &=  \left(\frac{1}{4 n_0 \alpha} W_{x}^2 -\frac{1}{n_0 \alpha} \frac{F_{x}^2}{W_{x}^2} - \frac{\gamma^2 W_{x}}{ 2 \alpha\, A^2}  \left<n(\vect{r}_{\perp},z), \pder{|u|^2}{W_{x}}\right>\right)  \\
    &\quad\left( \frac{1}{4 n_0 \alpha} W_{y}^2 - \frac{1}{n_0 \alpha} \frac{F_{y}^2}{W_{y}^2} - \frac{\gamma^2 W_{y}}{ 2 \alpha\, A^2}  \left<n(\vect{r}_{\perp},z), \pder{|u|^2}{W_{y}}\right>\right)\\
    &\quad + \frac{1}{4 n_0 \alpha}  \left(W_x^2+ W_y^2\right) \\    
    &\quad + \frac{1}{n_0 \alpha}  \left(\frac{F_{x}^2}{W_{x}^2} + \frac{F_{y}^2}{W_{y}^2}\right) \\
    &\quad -  \frac{\gamma^2}{2 \alpha A} \left< n(\vect{r}_{\perp},z),\pder{|u|^2}{A}\right>   \\ 
\end{split}
\end{equation}


\begin{equation}
\begin{split}
    \dot{P}  &= \frac{1}{2 n_0 \alpha} \left(W_{x}^2 + W_{y}^2\right)  - \frac{\gamma^2 W_{x}}{ 2 \alpha A^2}  \left<n(\vect{r}_{\perp},z), \pder{|u|^2}{W_{x}}\right>  
    - \frac{\gamma^2 W_{y}}{ 2 \alpha A^2}  \left<n(\vect{r}_{\perp},z), \pder{|u|^2}{W_{y}}\right>\\
    &\quad -  \frac{\gamma^2}{2 \alpha A} \left< n(\vect{r}_{\perp},z),\pder{|u|^2}{A}\right> 
\end{split}
\end{equation}


\begin{equation}
\begin{split}
    \dot{P}  &= \frac{1}{2 n_0 \alpha} \left(W_{x}^2 + W_{y}^2\right)  - \frac{\gamma^2 }{ 2 \alpha }\left<n(\vect{r}_{\perp},z), \frac{W_{x}}{A^2} \pder{|u|^2}{W_{x}}+ \frac{ W_{y}}{A^2}\pder{|u|^2}{W_{y}}+\frac{1}{A}\pder{|u|^2}{A}\right> 
\end{split}
\end{equation}



\pagebreak
\subsection{Equations}
\begin{equation}
\begin{split}
  A(z) &= A(0)\\
\dot{P}  &= \frac{1}{2 n_0 \alpha} \left(W_{x}^2 + W_{y}^2\right)  - \frac{\gamma^2 }{\alpha } \left<n(\vect{r}_{\perp},z), \frac{W_{x}}{2A^2} \pder{|u|^2}{W_{x}}+ \frac{ W_{y}}{2A^2}\pder{|u|^2}{W_{y}}+\frac{1}{2A}\pder{|u|^2}{A}\right> \\
\end{split}
\end{equation}

\begin{equation}
\begin{split}
 \dot{X} &= - 2T_x \\
\dot{Y} &= - 2T_y \\
\dot{T_x} &= - n_0 \gamma^2 \left<n(\vect{r}_{\perp},z), \frac{2}{A^2}\pder{|u|^2}{X}\right> \\
\dot{T_y} &= - n_0 \gamma^2 \left<n(\vect{r}_{\perp},z), \frac{2}{A^2} \pder{|u|^2}{Y}\right> \\
\end{split}
\end{equation}


\begin{equation}
\begin{split}
\dot{W_{x}} &= 4 F_{x}W_{x} \\
\dot{W_{y}} &= 4 F_{y}W_{y} \\
\dot{F}_{x}  &=   - \frac{1}{2 n_0 \alpha} W_{x}^4+ \frac{2}{n_0 \alpha} F_{x}^2  
+ \frac{\gamma^2}{\alpha}  \left<n(\vect{r}_{\perp},z), \frac{W_{x}^3}{A^2} \pder{|u|^2}{W_{x}}\right> \\    
\dot{F}_{y}  &=   - \frac{1}{2 n_0 \alpha} W_{y}^4 + \frac{2}{n_0 \alpha} F_{y}^2  + \frac{\gamma^2}{\alpha}  \left<n(\vect{r}_{\perp},z), \frac{W_{y}^3}{A^2} \pder{|u|^2}{W_{y}}\right>\\
\end{split}
\end{equation}


\subsection{Derivatives}

\begin{equation}
\begin{split}
|u|^2 = \frac{A^2W_x W_y}{\pi} \exp{-\left( W_x^2(x-X)^2 + W_y^2(y-Y)^2 \right)}
\end{split}
\end{equation}

\begin{equation}
\begin{split}
\pder{|u|^2}{A}  = 2A \frac{W_x W_y}{\pi} \exp{-\left( W_x^2(x-X)^2 + W_y^2(y-Y)^2 \right)}
\end{split}
\end{equation}

\begin{equation}
\begin{split}
\pder{|u|^2}{X}  &= \frac{A^2W_x W_y}{\pi} (2W_x^2(x-X)) \exp{-\left( W_x^2(x-X)^2 + W_y^2(y-Y)^2 
\right)}\\
\pder{|u|^2}{Y}  &= \frac{A^2W_x W_y}{\pi} (2W_y^2(y-Y)) \exp{-\left( W_x^2(x-X)^2 + W_y^2(y-Y)^2 \right)}
\end{split}
\end{equation}





\begin{equation}
\begin{split}
\pder{|u|^2}{W_{x}} &= \left( \frac{A^2 W_y}{\pi} - \frac{A^2 W_{x} W_y}{\pi} \left(2W_x(x-X)^2\right)  \right)\exp{-\left( W_x^2(x-X)^2 + W_y^2(y-Y)^2 \right)}\\
&= \left( 1 - W_{x} \left(2W_x(x-X)^2\right)  \right) \frac{A^2 W_y}{\pi} \exp{-\left( W_x^2(x-X)^2 + W_y^2(y-Y)^2 \right)}\\
&= \left(1  - 2 W_{x}^2 (x-X)^2  \right) \frac{A^2 W_y}{\pi}\exp{-\left( W_x^2(x-X)^2 + W_y^2(y-Y)^2 \right)}\\
\end{split}
\end{equation}


\begin{equation}
\begin{split}
\pder{|u|^2}{W_{y}} &= \left( \frac{A^2 W_x}{\pi} - \frac{A^2 W_{y} W_x}{\pi} \left(2W_y(y-Y)^2\right) \right)\exp{-\left( W_x^2(x-X)^2 + W_y^2(y-Y)^2 \right)}\\
&= \left( 1 - W_{y} \left(2W_y(y-Y)^2\right)  \right) \frac{A^2 W_x}{\pi} \exp{-\left( W_x^2(x-X)^2 + W_y^2(y-Y)^2 \right)}\\
&= \left(1  - 2 W_{y}^2 (y-Y)^2  \right) \frac{A^2 W_x}{\pi}\exp{-\left( W_x^2(x-X)^2 + W_y^2(y-Y)^2 \right)}\\
\end{split}
\end{equation}

\begin{equation}
\begin{split}
V = \frac{1}{2A}\pder{|u|^2}{A}  &= \frac{W_x W_y}{\pi} \exp{-\left( W_x^2(x-X)^2 + W_y^2(y-Y)^2 \right)}
\end{split}
\end{equation}

\begin{equation}
\begin{split}
\frac{W_{x}}{2A^2} \pder{|u|^2}{W_{x}} &= \frac{1}{2}\left(1  - 2 W_{x}^2 (x-X)^2  \right) \frac{W_{x} W_y}{\pi}\exp{-\left( W_x^2(x-X)^2 + W_y^2(y-Y)^2 \right)} \\
&= \frac{1}{2}\left(1  - 2 W_{x}^2 (x-X)^2  \right) V \\
\frac{W_{y}}{2A^2} \pder{|u|^2}{W_{y}} &= \frac{1}{2} \left(1  - 2 W_{y}^2 (y-Y)^2  \right) \frac{W_{x} W_y}{\pi}\exp{-\left( W_x^2(x-X)^2 + W_y^2(y-Y)^2 \right)}\\
&= \frac{1}{2} \left(1  - 2 W_{y}^2 (y-Y)^2  \right)  V \\
\end{split}
\end{equation}


\begin{equation}
\begin{split}
\frac{2}{A^2} \pder{|u|^2}{X}  &= \frac{W_x W_y}{\pi} (4W_x^2(x-X)) \exp{-\left( W_x^2(x-X)^2 + W_y^2(y-Y)^2 
\right)}\\
&=  4 W_x^2(x-X) V\\
\frac{2}{A^2} \pder{|u|^2}{Y}  &= \frac{W_x W_y}{\pi} (4W_y^2(y-Y)) \exp{-\left( W_x^2(x-X)^2 + W_y^2(y-Y)^2 \right)} \\
&=  4 W_y^2(y-Y) V\\
\end{split}
\end{equation}

\begin{equation}
\begin{split}
\frac{W_{x}^3}{A^2}\pder{|u|^2}{W_{x}} &= \left(1  - 2 W_{x}^2 (x-X)^2  \right) \frac{W_{x}^3 W_y}{\pi}\exp{-\left( W_x^2(x-X)^2 + W_y^2(y-Y)^2 \right)}\\
&= W_{x}^2 \left(1  - 2 W_{x}^2 (x-X)^2  \right) V \\
\frac{W_{y}^3}{A^2}\pder{|u|^2}{W_{y}} &= \left(1  - 2 W_{y}^2 (y-Y)^2  \right) \frac{W_{y}^3 W_x}{\pi}\exp{-\left( W_x^2(x-X)^2 + W_y^2(y-Y)^2 \right)}\\
&= W_{y}^2 \left(1  - 2 W_{y}^2 (y-Y)^2  \right) V \\
\end{split}
\end{equation}





\section{Calculating Statistics for Noise Terms: General}

\begin{equation}
\begin{split}
  A(z) &= A(0)\\
\dot{P}  &= \frac{1}{2 n_0 \alpha} \left(W_{x}^2 + W_{y}^2\right)  
- \frac{\gamma^2 }{\alpha } \left<n(\vect{r}_{\perp},z), V_{P}\right> \\
\end{split}
\end{equation}

\begin{equation}
\begin{split}
 \dot{X} &= - 2T_x \\
\dot{Y} &= - 2T_y 
\end{split}
\end{equation}

\begin{equation}
\begin{split}
\dot{T_x} &= - n_0 \gamma^2 \left<n(\vect{r}_{\perp},z),V_{T_{x}}\right> \\
\dot{T_y} &= - n_0 \gamma^2 \left<n(\vect{r}_{\perp},z),V_{T_{y}}\right>
\end{split}
\end{equation}


\begin{equation}
\begin{split}
\dot{W_{x}} &= 4 F_{x}W_{x} \\
\dot{W_{y}} &= 4 F_{y}W_{y}
\end{split}
\end{equation}

\begin{equation}
\begin{split}
\dot{F}_{x}  &=   - \frac{1}{2 n_0 \alpha} W_{x}^4+ \frac{2}{n_0 \alpha} F_{x}^2  + \frac{\gamma^2}{\alpha}  \left<n(\vect{r}_{\perp},z), V_{F_{x}}\right> \\    
\dot{F}_{y}  &=   - \frac{1}{2 n_0 \alpha} W_{y}^4 + \frac{2}{n_0 \alpha} F_{y}^2  + \frac{\gamma^2}{\alpha}  \left<n(\vect{r}_{\perp},z), V_{F_{y}}\right>
\end{split}
\end{equation}


\subsection{Modes}


\begin{equation}
\begin{split}
V = \frac{1}{2A}\pder{|u|^2}{A}  &= \frac{W_x W_y}{\pi} \exp{-\left( W_x^2(x-X)^2 + W_y^2(y-Y)^2 \right)}
\end{split}
\end{equation}

\begin{equation}
\begin{split}
V_{P} &= \frac{W_{x}}{2A^2} \pder{|u|^2}{W_{x}} + \frac{ W_{y}}{2A^2}\pder{|u|^2}{W_{y}} + \frac{1}{2A}\pder{|u|^2}{A} \\
& =  \frac{1}{2}\left(\left(1  - 2 W_{x}^2 (x-X)^2  \right)  +  \frac{1}{2} \left(1  - 2 W_{y}^2 (y-Y)^2  \right)  + 1\right) V \\
& =  \left(2  - W_{x}^2 (x-X)^2  - W_{y}^2 (y-Y)^2  \right) V 
\end{split}
\end{equation}


\begin{equation}
\begin{split}
V_{T_{x}} =  4 W_x^2 (x-X) V
\end{split}
\end{equation}

\begin{equation}
\begin{split}
V_{T_{y}} =  4 W_y^2 (y-Y) V
\end{split}
\end{equation}

\begin{equation}
\begin{split}
V_{F_{x}} = W_{x}^2 \left(1  - 2 W_{x}^2 (x-X)^2  \right) V \\
\end{split}
\end{equation}

\begin{equation}
\begin{split}
V_{F_{y}} = W_{y}^2 \left(1  - 2 W_{y}^2 (y-Y)^2  \right) V \\
\end{split}
\end{equation}





It is natural to define the innerproduct as
\begin{equation}
\begin{split}
\ip{f}{g} = \Ip{f}{g} 
\end{split}
\end{equation}
and then terms of the form 
\begin{equation}
\begin{split}
\ip{n(\vect{r},z)}{V(\vect{r},\vect{p}(z))} = \Ip{n(\vect{r},z)}{V(\vect{r},\vect{p}(z))}
\end{split}
\end{equation}
will have statistics 
\begin{equation}
\begin{split}
\Expect{\ip{n(\vect{r},z)}{V(\vect{r},\vect{p}(z))}} = \Ip{\Expect{n(\vect{r},z)}}{V(\vect{r},\vect{p}(z))}
\end{split}
\end{equation}

\begin{equation}
\begin{split}
\Expect{\ip{n(\vect{r},z_1)}{V(\vect{r},\vect{p}(z_1))}\ip{n(\vect{r},z_2)}{V(\vect{r},\vect{p}(z_2))}} 
= \iint \iint   \Expect{n(\vect{r}_1,z_1)n(\vect{r}_2,z_2)} V(\vect{r}_1,\vect{p}(z_1))V(\vect{r}_2,\vect{p}(z_2)) \,d\vect{r}_1\,d\vect{r}_2
\end{split}
\end{equation}



\subsection{Markov Approximation}
If we assume the index is delta correlated in the propagation direction, than the correlations are easier to work with.
\begin{equation}
\begin{split}
 \Expect{n(\vect{r}_1,z_1)n(\vect{r}_2,z_2)} = C(|\vect{r}_1-\vect{r}_2|)\delta(z_1-z_2) = C(x_1-x_2,y_1-y_2)\delta(z_1-z_2)
\end{split}
\end{equation}

\subsubsection{P}
\begin{equation}
\begin{split}
&\Expect{\ip{n(\vect{r}_1,z_1)}{V_{P}(\vect{r}_1;W_x(z_1),W_y(z_1))}\ip{n(\vect{r},z_2)}{V(\vect{r}_2;W_x(z_2),W_y(z_2))}} \\
&= \iint \iint   C(x_1-x_2,y_1-y_2) V_{P}(x_1,y_1;W_x(z_1),W_y(z_1)) V_{P}(x_2,y_2;W_x(z_1),W_y(z_1)) \,d\vect{r}_1\,d\vect{r}_2 \delta(z_1-z_2)
\end{split}
\end{equation}



%\begin{align}&\left< \left[\begin{array}{c}
%F_{E}(z) \\ 
%F_{W}(z) \\
%F_{T}(z) \\
%F_{C}(z) \\ 
%F_{\Omega}(z)\\ 
%F_{\phi}(z) \\
%\end{array}\right] \left[\begin{array}{c}
%F_{E}(z') \\ 
%F_{W}(z') \\
%F_{T}(z') \\
%F_{C}(z') \\
%F_{\Omega}(z')\\ 
%F_{\phi}(z') \\
%\end{array}\right]^{T} \right>=\\
%& \frac{\sigma^2}{\sqrt{\pi}} \left[\begin{array}{c|cccccc}
%           & F_{E} & F_{W}           & F_{T}                &  F_{C}              & F_{\Omega}                                         & F_{\phi}                                               \\
%\hline                                                                                                                                                                                          \\
%F_{E}      & 2E    & 0               & 0                    &  0                  & 0                                                  & 0                                                      \\
%F_{W}      & 0     & \frac{W^2}{ E } & 0                    &  0                  & 0                                                  & 0                                                      \\
%F_{T}      & 0     & 0               & \frac{W^2}{E}        &  0                  & \frac{ CW^2 }{ E}                                  & \frac{\Omega W^2 }{ E}                                 \\ 
%F_{C}      & 0     & 0               & 0                    &  \frac{4}{ E W^4 }  & 0                                                  & -\frac{1}{2 E W^2}                                      \\
%F_{\Omega} & 0     & 0               & \frac{CW^2}{E}       &  0                  & \frac{1}{E} \left( \frac{1}{W^2} + C^2W^2 \right)  & \frac{C\Omega  W^2 }{ E }                              \\
%F_{\phi}   & 0     & 0               & \frac{\Omega W^2}{E} &  -\frac{1}{2 E W^2}    & \frac{C\Omega W^2 }{E }                            & \frac{1}{E} \left(\frac{3}{4}   + \Omega^2 W^2 \right) \\ 
%\end{array}\right] \delta(z-z')
% \end{align}

 \subsection{Rytov Expansion}
 Look for solutions of the form
 \begin{equation}
   A = \exp{B} \quad \mbox{which gives} \quad \nabla A = \left(\nabla B\right) \exp{B} \quad
 \mbox{and} \quad \nabla^2 A = \left(\left|\nabla B \right|^2 + \nabla^2 B\right)\exp{B}
 \end{equation}
 and
 \begin{equation}\label{EQ:helmholtz_equation_3}
 \begin{split}
   &\left|\nabla B \right|^2 + \nabla^2 B + k^2  (n_0 + \epsilon n_1(x,y)) = 0 \quad  \mbox{where} \quad B(x,y,0) = B_{i}(x,y) + i \epsilon_{1} \Phi(x,y).
 \end{split}
 \end{equation}
 Now expand $B$ in an expansion of the form
 \begin{equation}
 \begin{split}
 B = B_{0} + \epsilon B_{1} + \dots
 \end{split}
 \end{equation}

 % At $\bigOh{1}$:
 \begin{equation}\label{EQ:Order_1}
 \begin{split}
   &\left|\nabla B_{0} \right|^2 + \nabla^2 B_{0} + k^2 n_0 = 0 \quad  \mbox{where} \quad B(x,y,0) = B_{i}(x,y).
 \end{split}
 \end{equation}
 and letting $B_{0} = \ln{A_{0}}$ gives
 \begin{equation}
 \begin{split}
   \nabla^2 A_{0} + k^2n_0 A_{0} = 0  \quad  \mbox{where} \quad A_{0}(x,y,0) = A_{i}(x,y),
 \end{split}
 \end{equation}
 the solution of which is given by the Green's function for the Helmholtz Equation,
 \begin{equation}
 \begin{split}
 A(x,y,z;\epsilon_1,\epsilon_2) = - i  \frac{\epsilon_2}{\epsilon_1} \exp{i  \frac{2\pi}{\epsilon_1}}
 \iint A_{i}(x,y) \exp{i \pi  \frac{\epsilon_2}{\epsilon_1} \left[ \left(x-x^*\right)^2 + \left(y-y^*\right)^2  \right]} \, dx \, dy
 \end{split}
 \end{equation}



%---------------------------------------------------
%--------------biblatex stuff --------------------
%---------------------------------------------------
%  

%\bibliographystyle{abstract}
 \bibliographystyle{plain}
 \bibliography{./Bibliography/Bibliography}

% \thispagestyle{empty}
% \printbibliography

%---------------------------------------------------
%----------------------LAYOUT --------------------
%---------------------------------------------------
\label{lastpage}
\pagebreak
\end{document}
